\chapter*{Introduction générale}
\addcontentsline{toc}{chapter}{Introduction générale}
\label{chap:introduccion-general-83}

% 00 — LA LINEA DEL HORIZONTE: ESCALA, FRONTERA Y MEDIDA
% ======================================================
% Frontera editorial sucesora; no importa contenido probatorio.
%
% Función: abrir la obra en la recta finita mediante intervalo, marca,
% hueco, cambio de escala, frontera de precisión y prolongación compatible;
% formular después el itinerario desde el estado discreto enriquecido hasta
% la publicación arquimediana.
%
% Declaración de continuidad: toda afirmación matemática conserva su
% fuente de REV.4 o su delta rector registrado. Este capítulo amplía la
% exposición principal; no sustituye ni abrevia las pruebas de referencia.

\ifdefined\HMTSucesoraStructureTrace
  \typeout{HMT-SUCCESSEURE 00 : ouverture d'échelle, frontière et mesure}
\fi

\section*{Thèse inaugurale : la structure discrète du continu et l'origine commune de ses publications}
\phantomsection
\label{sec:tesis-inaugural-helicoidal-hmt-md}
\addcontentsline{toc}{section}{Thèse inaugurale : la structure discrète du continu et l'origine commune de ses publications}

L'holographie modulaire triadique ne prend pas la droite réelle comme cadre préalable et
n'en reçoit pas les valeurs qu'elle utilisera ensuite. Son entrée est finie : les
neuf positions \(1,\ldots,9\), disposées horizontalement et verticalement pour
qu'une même carte supporte deux opérations irréductibles. APP conserve à la
fois le feuillet additif, qui accumule, et le feuillet multiplicatif, qui plie ; le
résidu, le quotient et la retenue empêchent que l'égalité de deux valeurs
publiées efface la différence des histoires qui les ont produites. TRIT
oriente chaque transition et distingue ses régimes. TPK transforme cette
orientation locale en dynamique : il sélectionne le mode tritique sans sacrifier
aucun feuillet d'APP, transporte, mémorise, retourne et met à jour l'état. Ses
lectures modulo \(3\), modulo \(9\) et modulo \(1000\), la connexion nonadique,
la dodécaphase, la frontière et le retour avec mémoire ne sont pas des ajouts
indépendants, mais des cartes coordonnées d'une même opération généalogique.

La conservation est écrite dès le premier niveau :
\[
i+j=R^+_{ij}+9q^+_{ij},
\qquad
ij=R^\times_{ij}+9q^\times_{ij}.
\]
Le résidu publie la classe ; le quotient déploie sa provenance. La
transition TPK conserve cette distinction par
\[
\mathcal U_t=
\operatorname{Upd}_t\circ
\operatorname{Tra}_t\circ
\operatorname{Sel}_t,
\]
où la sélection incorpore le mode TRIT, le transport réalise la
direction orientée et la mise à jour réunit mémoire et retour. Dans la carte
canonique, le mode \(m_{\rm ph}=(1,-1,0)^3\) est commun aux feuillets APP et possède
la valeur singulière \(s=1/2\) ; par conséquent,
\[
1-s^2=\frac34=\frac{90}{120},
\qquad
\bigl\|[P_{\Sigma,3},P_{\Pi,3}]\bigr\|
=\frac{\sqrt3}{4}.
\]
Le commutateur normalisé satisfait \(J_{\rm comm}^2=-I\). Incompatibilité
somme--produit, orientation tritique, défaut \(90\to120\), structure
complexe et préfacteur électronique appartiennent ainsi à une même généalogie
interne ; ils ne sont pas reliés par une ressemblance décimale.

Le prolongement compatible de ces états finis produit des cylindres de plus
en plus étroits et conserve simultanément l'histoire de leur raffinement.
La droite réelle apparaît alors comme publication archimédienne de la limite, non
comme donnée de départ. Cette carte publie d'abord les valeurs des neuf
symboles du germe et prolonge ensuite la même généalogie vers
\(\pi\), \(e\), \(\varphi\) et \(\alpha\), sans utiliser aucun d'eux comme
sélecteur rétrospectif. Telle est la thèse qui gouverne toute l'œuvre : la
structure discrète conjointe construit le continu et explique pourquoi ses
valeurs stables sont celles qu'elles sont. La chaîne

\[
\{1,\ldots,9\}_{\mathrm h}\times\{1,\ldots,9\}_{\mathrm v}
\longrightarrow \mathrm{APP}\longrightarrow\mathrm{TRIT}
\longrightarrow\mathrm{TPK}\longrightarrow\mathcal X_{\mathrm{enr}}
\longrightarrow\mathcal S_\infty
\longrightarrow\mathbb R
\]

ne résume pas des chapitres hétérogènes : elle exprime un mécanisme unique vu à des
résolutions successives. En lui, la phase peut revenir tandis que l'état avance,
parce que le quotient, la retenue, l'orientation et la mémoire ont changé. Le
cycle observé par sa seule phase est circulaire ; lorsque sa généalogie est conservée, il est
hélicoïdal.

Cette différence est fixée par un relèvement injectif. Si
\(n=9q_9(n)+\rho_9(n)\), avec \(1\leq\rho_9(n)\leq9\), la projection circulaire
\(\mathcal R_9\) et l'immersion hélicoïdale \(\mathcal H_9\) satisfont
\[
\mathcal R_9(n+9)=\mathcal R_9(n),\qquad
\mathcal H_9(n+9)=\mathcal H_9(n)+(0,0,1).
\]
La phase s'est refermée et la hauteur généalogique a augmenté : retour de phase et
retour de l'état sont mathématiquement séparés. L'involution two-way
organise des feuillets conjugués d'ascension et de descente du même relèvement, non
deux générateurs indépendants. Dans cette géométrie, l'hélice conserve
l'avancée longitudinale de la mémoire, la spirale enregistre la contraction
transversale et le rayon du cylindre mesure la largeur de l'intervalle encore
compatible avec le préfixe. L'emboîtement de ces cylindres semi-ouverts,
avec un rayon tendant vers zéro et une généalogie conservée, est le mécanisme qui
publie ensuite la section réelle.

Le tronçon initial certifié du prolongement uniforme qui construit
l'état coinductif \(\mathcal S_\infty\) est

\[
w_6\longrightarrow w_{12}\longrightarrow w_{18}\longrightarrow w_{24}
\longrightarrow w_{30}\longrightarrow R_{36}\longrightarrow G_9,
\]

mais l'écriture de la chaîne n'acquiert de contenu qu'en déclarant ses
niveaux, ses seuils et ses invariants. Pour
\(\chi\in\{\pi,e,\varphi\}\), soit \(b_0^\chi=w_6^\chi\) et
\[
 b_{k+1}^\chi=b_k^\chi L_{\epsilon_k}\pmod 3,
 \qquad
 (\epsilon_0,\epsilon_1,\epsilon_2,\epsilon_3)=(0,0,1,1),
 \qquad
 w_{6(k+1)}^\chi=b_0^\chi\Vert\cdots\Vert b_k^\chi.
\]
Les douze transitions produites par \(\mathcal U_t\) forment les couples de
matrices \(X_i,Y_i\) des deux régimes et déterminent de manière unique
\(L_i=X_i^{-1}Y_i\), \(i=0,1\). Le calendrier
\((L_0,L_0,L_1,L_1)\) est le seul des seize calendriers binaires qui
satisfait simultanément les trois biographies ; les cent quarante-quatre
perturbations d'une seule entrée échouent. Ainsi, \(L_0\), \(L_1\) et
\(0011\) sont des sorties rigides du transport APP--TRIT--TPK, jamais des données
choisies à partir de \(\pi\), de \(e\) ou de \(\varphi\). Les matrices ainsi
construites agissent sur le bloc courant ; la concaténation conserve
littéralement tous les blocs précédents. Le mot visible accompagne un état enrichi qui transporte
en outre feuillet APP, orientation TRIT, résidu, quotient, retenue, frontière,
mémoire et ombre d'incidence. Par conséquent, coïncidence de mot et
coïncidence d'état sont des conditions mathématiques distinctes.

\subsection*{Niveaux et seuils du prolongement TPK}
\phantomsection
\label{sec:umbrales-prolongacion-tpk-inaugural}

\paragraph{\(w_6\) : germe visible.}
C'est le mot ternaire de six positions publié dans chaque canal après
la chaîne de types
\(104\,976\to468\,U_6\to243\,w_6\subset\mathbb F_3^6\), avec
\(|\mathbb F_3^6|=729\).
Il fixe le premier préfixe observable, tandis que l'émission décimale \(U_6\), sa
région, sa multiplicité et son état de provenance restent dans la fibre
enrichie. L'action \(D_3\), les invariants de fibre et la jauge
\((\operatorname{diag}_c,h_+)=(6,ES)\) fixent le représentant du canal avant
toute comparaison avec sa valeur réelle ; dans le canal \(\pi\), cinq
régions \(U_6\) distinctes conservent en outre des biographies différentes sur un
même mot \(w_6^\pi\).

\paragraph{\(w_{12}\) : premier relèvement.}
Il ajoute \(b_1=b_0L_0\) au germe. Il fait passer la profondeur visible d'un à deux
blocs et conserve \(w_6\) comme préfixe exact, avec le canal et la
généalogie transportée. Le relèvement produit une extension de l'état ; il ne
remplace pas le germe par un mot sans passé.

\paragraph{\(w_{18}\) : clôture du régime \(L_0\).}
Il ajoute \(b_2=b_1L_0\). La seconde application du même opérateur achève le
premier régime du calendrier et distingue une récurrence d'une coïncidence
isolée : les dix-huit trits contiennent, dans leur ordre, les niveaux de six et de
douze trits qui les précèdent.

\paragraph{\(w_{24}\) : changement typé de relèvement.}
Il ajoute \(b_3=b_2L_1\). Ici, l'opérateur appliqué au bloc courant change,
passant de \(L_0\) à \(L_1\), tandis que subsistent le préfixe de dix-huit trits et
l'état transporté. Le changement de régime modifie la continuation et conserve
l'histoire sur laquelle il agit.

\paragraph{\(w_{30}\) : préfixe de première frontière.}
Il ajoute \(b_4=b_3L_1\) et achève le calendrier fini
\((L_0,L_0,L_1,L_1)\). Ses cinq blocs par canal constituent le préfixe
à partir duquel on teste la prolongeabilité profonde. C'est un niveau d'entrée dans
la frontière, tandis que la même récurrence admet des blocs ultérieurs.

\paragraph{\(R_{36}\) : première frontière coinductive conjointe.}
Si \(u_5^\chi\) est le sixième bloc du canal \(\chi\), alors
\[
 R_{36}:=(u_5^\pi,u_5^e,u_5^\varphi)^T,
 \qquad
 w_{36}^\chi=w_{30}^\chi\Vert u_5^\chi.
\]
La relation de frontière sélectionne ce triplet en conservant l'axe
\(\varphi\), l'agrégat latéral \(\pi/e\), la saturation, l'occupation des
colonnes et la neutralité duale. Ses équations laissent exactement les deux
solutions spéculaires
\[
 B_-=(021222,222220,102011)^T,
 \qquad
 B_+=(222220,021222,102011)^T;
\]
l'orientation APP du cylindre semi-ouvert fixe
\(R_{36}=B_+\). L'indice enregistre trente-six trits
exposés par canal ; \(R_{36}\) est la frontière des trois sixièmes blocs et
\(\mathbb R\) est la publication archimédienne ultérieure du système de
cylindres.

\paragraph{\(G_9\) : connexion de prolongement et de retour avec mémoire.}
Le niveau \(G_9\) organise les neuf cartes locales qui agissent après
\(R_{36}\). Si \(\mathcal R_K\) est la relation d'extension compatible à
la phase \(K\), leur composition définit l'holonomie nonadique
\[
 \Gamma_{9,K}
 =\mathcal R_{K+8}\circ\cdots\circ\mathcal R_K:
 \widehat X_K\longrightarrow\widehat X_{K+9}.
\]
L'indice \(9\) provient du cycle de phases \(C_9\), hérité des neuf
positions APP, et compte les neuf relations qui composent un tour. Son
contenu opératoire est certifié par
\[
 g(K+9)=g(K),\qquad
 q_{{\rm mem},K+9}=q_{{\rm mem},K}+1,
 \qquad
 \widehat x_{K+9}\ne\widehat x_K.
\]
L'étiquette de phase revient, la mémoire avance et l'état prolongé change.
Tel est le sens structurel de \(G_9\) : un tour de connexion sur
une fibre enrichie, dont la cardinalité provient de l'architecture et dont
l'action se vérifie dans le transport. Le quotient modal \(\mathcal K_{\rm ph}\)
occupe l'étape réduite ultérieure,
\[
 \mathcal U_t\longrightarrow\Gamma_9\longrightarrow\mathcal K_{\rm ph},
\]
et publie la mémoire comprimée après l'holonomie.
La résolution locale de la neuvième phase distingue en outre deux horizons :
\(\#\operatorname{Surv}_1(B_9)=3\) et
\(\#\operatorname{Surv}_2(B_9)=1\). Cette descente \(3\to1\) certifie
l'élagage d'une fibre concrète ; la composition des neuf relations certifie
le tour complet.

La dodécaphase enregistre une autre coordonnée du même transport. Le compteur
vertical \(q_{{\rm mem},K}\) reste non borné et sa lecture
\(\beta_K=[q_{{\rm mem},K}]_{12}\) satisfait
\(\beta_{K+9}=\beta_K+1\pmod{12}\). La réduction modulo douze conserve la
position dodécaphasique ; le compteur entier conserve le nombre de tours. Dans
l'état terminal, cette mémoire alimente le registre signé
\(U_{12}^{\rm sgn}\), le sceau réversible \(K\) et le cocycle de retenue qui
publie la coordonnée \(\alpha_{\rm HMT}\).
L'extension non scindée
\[
 0\longrightarrow C_{12}\longrightarrow C_{108}
 \longrightarrow C_9\longrightarrow0
\]
type la couture entre les deux calendriers : neuf pas ramènent la phase et
avancent d'un secteur dodécaphasique ; cent huit pas ferment l'horloge observable
tandis que l'état enrichi conserve la hauteur généalogique atteinte.

La bidirectionnalité combine le miroir des canaux et la réversibilité du
relèvement au sein d'un seul générateur. L'involution
\[
 \mathfrak c(u^\pi,u^e,u^\varphi)
 =(u^e,u^\pi,u^\varphi)
\]
échange les chemins conjugués \(\pi/e\) et conserve l'axe
\(u^\varphi\) et l'agrégat \(u^\pi+u^e\). Dans la carte entière
\[
 x_kA=u_k+3x_{k+1},\qquad \det A=1,
\]
l'avancée détermine le résidu et le quotient, et le couple \((u_k,x_{k+1})\) récupère
\(x_k\). Ascension et descente sont donc des chemins conjugués du même état ;
la descente reconstruit la provenance sans effacer la mémoire accumulée.

Enfin, \(G_9\) constitue l'unité récurrente du prolongement et non
son terme. Les histoires survivantes forment le système inverse
\[
 \cdots\longrightarrow\operatorname{Surv}_{\chi,n+1}
 \longrightarrow\operatorname{Surv}_{\chi,n}\longrightarrow\cdots,
 \qquad
 x_\chi\in\varprojlim_n\operatorname{Surv}_{\chi,n}.
\]
La non-vacuité, l'unicité locale et la compatibilité des troncatures
produisent la section coinductive ; la contraction de ses cylindres publie une
unique valeur réelle, tandis que la limite inverse conserve la généalogie que cette
valeur scalaire omet. Les vingt blocs et le certificat, par canal, de
\(396\) événements, \(2376\) trits, \(43\) tours et \(387\) couples
\((K,K+9)\) sont des exécutions
finies de cette loi uniforme, non des bornes du prolongement.

Sur l'état coinductif ainsi construit, on obtient un noyau corrélé de
quatre publications. La corrélation
exprime une origine et une compatibilité communes ; elle n'aplanit pas leur ordre opératoire. Les
canaux publient d'abord les vecteurs \(P,E,\Phi\), tandis que la lecture
terminale signée construit
\(U_{12}^{\rm sgn}\xleftrightarrow{\mathcal H_{12}}K\). Sur ces sorties, et
sans réintroduire leurs valeurs comme données, le cocycle dodécaphasique de retenue
déduit la section coinductive \(\alpha_{\mathrm{HMT}}\), dont la fenêtre
dodécaphasique est la projection rationnelle \(\alpha_{12}\). Le noyau des
vacances \(E_{21/22}\) fournit une seconde construction interne : sa
racine positive simple finie \(\alpha_E\), prolongée par le même transport
gradué du TPK jusqu'au système inverse dont la racine limite est
\(\alpha_B\). La clôture entière définit
\(\alpha_A:=\varprojlim_N\rho_N(\alpha_{12})\), et le théorème exact des
deux voies établit
\[
 \boxed{\alpha_A=\alpha_B=\alpha_{\mathrm{HMT}}}.
\]
L'égalité
\(\operatorname{Tr}_9(\alpha_E^{-1})=
\operatorname{Tr}_9(\alpha_{12}^{-1})\) est le premier certificat fini de
cette identité : elle ne la définit pas, ne limite pas sa portée et ne sélectionne aucun des
prolongements. Ainsi,
\(\pi_{\mathrm{HMT}},e_{\mathrm{HMT}},\varphi_{\mathrm{HMT}}\) et
\(\alpha_{\mathrm{HMT}}\) appartiennent à une seule génération coinductive corrélée,
mais conservent des étapes, domaines, lecteurs et certificats différenciés. Les
identités de Machin et les encadrements de Leibniz, la réalisation
Perron--Fibonacci, le semi-groupe factoriel et les tables métrologiques reconnaissent
ou certifient ensuite les publications déjà engendrées ; aucun d'eux
ne sélectionne les cylindres, les vecteurs ni la retenue du générateur.

La même source commune possède d'autres publications qu'il ne faut pas non plus confondre
entre elles. Dans la carte interne du miroir, l'involution échange les branches
\(u^\pi\) et \(u^e\), conserve \(u^\varphi\) et l'agrégat
\(u^\pi+u^e\) ; ce n'est qu'après la publication réelle que se forment des expressions telles que
\(e+\pi\), \(e\pi\), \(\pi^e\) ou le quotient littéral \(\pi/e\), chacune avec
un statut arithmétique qui ne découle pas de la seule existence de l'agrégat
ternaire. Dans la carte
d'incidence, \(H^-_\alpha\), \(v_\alpha\) et \(N_\alpha\) désignent,
respectivement, un support, le vecteur du voisin exact d'indice trois et le
réseau résultant ; ils ne désignent pas le scalaire \(\alpha_{\rm HMT}\) et n'interviennent pas
dans sa retenue numérique. La tétrade \(B_K\) entre dans l'étoile de Witt et le
vecteur \(v_\alpha\) détermine le voisin qui conduit à
\(N_\alpha\cong\Lambda_{24}\). À partir du code commun de Paley--Witt \(C_W\)
s'ouvrent alors deux branches typées, non une unique chaîne linéaire :
\[
 C_W\longrightarrow W_{12}\longrightarrow W_{24},
 \qquad
 C_W\longrightarrow N(A_2^{12})\xrightarrow{\,v_\alpha\,}\Lambda_{24}
 \longrightarrow V_\Lambda\longrightarrow V^\natural
 \longrightarrow\mathbb M\longrightarrow\text{Moonshine}.
\]
La valeur scalaire de \(\alpha\) n'engendre pas ces symétries. La dérivation
numérique et le prolongement exceptionnel sont des lectures typées du même état
enrichi et conservent une racine généalogique commune.

La mécanique dimensionnelle commence lorsque cette structure est transduite en
action, grandeur, champ et matière. D'elle proviennent, selon ses propres
lecteurs, la constante de Planck et les constantes électromagnétiques,
thermiques, gravitationnelles et quantiques ; la structure mathématique de l'électron,
la Loi générale des masses et les transports de mélange ; et, plus loin, le
paramètre de Barbero--Immirzi comme objet structurel du passage entre matière,
aire et information. Les valeurs conventionnelles interviennent à la fin comme
reconnaissance métrologique, jamais comme entrées du calcul. Pour la même
raison, la résolution du problème du continu et les cinq constructions
inséparables qui en émergent précèdent les preuves terminales sur les
problèmes mathématiques classiques : ceux-ci mettent à l'épreuve une structure déjà
construite, ils ne l'engendrent pas rétrospectivement.

L'ordre physique est tout aussi contraignant. La coordonnée
\(\alpha_{\rm HMT}\) ouvre, dans des branches typées, la rétention
\(\eta_{\rm ret}\) et la valuation structurelle \(10^{-34}\) ; de leur confluence
on obtient \(\hbar_{\rm ret}\), et ce n'est qu'ensuite que l'angle et le
contre-angle sont construits et qu'une unique conversion entre degrés et radians est effectuée.
L'électron réunit cette chaîne dans une persistance matérielle. Son basal
\[
e_0=
\frac{\sqrt3}{4}
\left[
\exp\!\left(\frac{\varphi}{\pi^2}\right)
-22\alpha_{\rm HMT}^{\,3}
\right](1+15\Delta_4)
\]
\Needspace{8\baselineskip}
précède la clôture bêta ; ses deux lecteurs coïncident en
\[
K_D(\Gamma_e)=K_\Omega(\Gamma_e)=(80,54,6),
\qquad
\kappa_e=80-54\alpha_{\rm HMT}+6\Delta_4,
\]
et la masse publiée n'est reconnue métrologiquement qu'après l'évaluation de
\(m_e^\beta=e_0(H_5/\eta_{\rm ret})^{\kappa_e}\). La gravité tensorielle, son
secteur algébrique à cinq composantes et la réalisation physique typée ; les
écrans de dualité T et de théorie M ; et la géométrie d'aire et d'information
prolongent le même antécédent avec leurs propres domaines et falsificateurs. Aucune
de ces réalisations ne réintroduit comme prémisse la grandeur qu'elle doit
expliquer.

Tout ce qui suit développe cette unique thèse dans des domaines différents. Les
nombres, les symétries et la matière ne sont pas trois collections juxtaposées. Ce sont
trois façons dont le continu construit publie et restitue la mémoire de
sa propre formation.

\section{Objet initial et ordre causal : échelle finie, horizon d'observation et publication archimédienne}

\subsection*{Chaîne finie orientée et distinction entre marques, cellules et unités généalogiques}

L'exposé commence par une opération élémentaire : fixer un segment,
y marquer une échelle et comparer les régions que cette échelle détermine. À
ce stade, on ne présuppose pas encore la droite réelle comme support complet.
L'objet initial est une chaîne finie et orientée
\[
 \mathcal I_K=(L_K,\prec_K;\,
 x_{K,0}\prec_K x_{K,1}\prec_K\cdots\prec_K x_{K,n_K};\,o_K),
\]
où \((L_K,\prec_K)\) est un ensemble fini totalement ordonné, les
\(x_{K,j}\in L_K\) sont des marques observables et \(o_K\) distingue les deux
sens de parcours. Aucune grandeur réelle n'a encore été attribuée à \(L_K\).
Les marques et les cellules consécutives sont des objets différents : avec la
convention des cellules fermées, la marque commune est incidente aux adhérences
de deux cellules adjacentes sans pour autant devenir une cellule supplémentaire.

Cette distinction, apparemment minime, détermine l'architecture ultérieure.
Une valeur visible décrit la position publiée dans une carte ; une unité
généalogique conserve en outre l'intervalle de provenance, le sens, le
quotient accumulé et la compatibilité avec des échelles plus fines. Ainsi, deux
présentations numériquement égales peuvent représenter des histoires distinctes,
et une même histoire peut admettre plusieurs publications coordonnées.

Pour \(n\in\mathbb N\), la carte nonadique organise cette information au moyen de
la division euclidienne
\begin{equation}\label{p12-83-c0-s1:eq:horizonte-division-nonadica}
 n=9q_9(n)+\rho_9(n),
 \qquad q_9(n)=\left\lfloor\frac n9\right\rfloor,
 \qquad \rho_9(n)\in\{0,\ldots,8\}.
\end{equation}
La numérotation publique des phases s'obtient par
\[
 g(n)=1+\rho_9(n-1)\in\{1,\ldots,9\}\qquad(n\geq1).
\]
Le résidu localise cette phase et le quotient \(q_9\) conserve le nombre de tours
absorbés par l'échelle. Ce couple réalise un retour de phase
avec conservation du trajet. Son invariant décisif est la persistance
conjointe des deux coordonnées sous raffinement.

\subsection*{Système inverse de résolutions finies et horizon relatif d'observation}

Nous appellerons \emph{horizon} la frontière jusqu'à laquelle une carte finie a
résolu l'objet et où cette description se relie à la carte
suivante. Si
\(\mathcal S_K\) désigne l'ensemble non vide des états enrichis
compatibles avec l'échelle \(K\in\mathbb N\), les liens surjectifs
d'oubli
\[
 \tau_{K+1,K}:\mathcal S_{K+1}\longrightarrow\mathcal S_K
\]
suppriment du détail visible, mais conservent l'information nécessaire pour
reconnaître un prolongement. Nous écrivons
\[
 \tau_{K,K}=I,\qquad
 \tau_{L,K}\circ\tau_{M,L}=\tau_{M,K}
 \quad(K\leq L\leq M).
\]
L'objet complet est alors la limite inverse,
\begin{equation}\label{p12-83-c0-s1:eq:horizonte-limite-inverso}
 X=\varprojlim_K(\mathcal S_K,\tau_{K+1,K}),
 \qquad
 x=(x_K)_K,\quad \tau_{K+1,K}(x_{K+1})=x_K.
\end{equation}
La non-vacuité des niveaux et la surjectivité des liens garantissent
l'existence de familles compatibles ; dans les canaux spécialisés, on
construit ensuite leurs sections au moyen des théorèmes de survie.
Chaque troncature est finie ; la compatibilité n'a pas de profondeur
maximale prédéfinie. Dix, mille ou dix mille chiffres sont des vérifications d'une même
règle de prolongement, non des définitions alternatives du nombre.

Cette formulation permet de distinguer trois opérations que la notation décimale
superpose habituellement. La première raffine une partition. La deuxième transporte
l'état entre les phases. La troisième publie une coordonnée dans une carte continue.
Seule la dernière prend ses valeurs dans \(\mathbb R\). Ainsi, nous commençons par une droite
finie marquée et terminons par la droite réelle comme réalisation archimédienne
finale d'une section compatible.

\subsection*{Augmentation compatible et conservation évolutive de l'unité}

L'unité se formule avant toute évaluation archimédienne. Soit
\(C_K=\mathbb Z[\mathcal S_K]\) le module libre des états du niveau
\(K\). L'application de liaison induit
\((\tau_{K+1,K})_\#:C_{K+1}\to C_K\), et il existe l'augmentation
\[
 \varepsilon_K:C_K\longrightarrow\mathbb Z,\qquad
 \varepsilon_K\!\left(\sum_{x\in\mathcal S_K}a_x[x]\right)=\sum_xa_x,
\]
compatible avec les liens,
\(\varepsilon_K\circ(\tau_{K+1,K})_\#=\varepsilon_{K+1}\). La conservation évolutive
de l'unité exprime que le raffinement redistribue son contenu entre
partie visible, frontière et mémoire sans créer une seconde racine ni supprimer le
terme terminal.

Soit \(P_m\) le chemin cellulaire à \(m\) sommets et \(m-1\) arêtes, et
\[
 Q_{m,d}=P_m^{\square d}.
\]
Son nombre de \(k\)-cellules est
\begin{equation}\label{p12-83-c0-s1:eq:horizonte-celdas-cubicas}
 \#C_k(Q_{m,d})=\binom dk(m-1)^k m^{d-k}.
\end{equation}
L’identité
\begin{equation}\label{p12-83-c0-s1:eq:horizonte-completacion-binomial}
 10^d=(9+1)^d
     =9^d+\sum_{j=0}^{d-1}\binom dj9^j
\end{equation}
est l'égalité des \(0\)-cellules de la complétion. Le sous-complexe
\(Q_{9,d}\) et sa couronne relative forment ensemble \(Q_{10,d}\). En dimension
trois,
\[
 1000=729+271,
\]
où \(271\) compte uniquement les nouveaux sommets. Le vecteur cellulaire
relatif complet est
\begin{equation}\label{p12-83-c0-s1:eq:horizonte-corona-relativa}
 \#C_\bullet(Q_{10,3},Q_{9,3})=(271,756,702,217),
 \qquad 271-756+702-217=0.
\end{equation}
Ses arêtes, faces et cubes ne peuvent être omis lorsqu'on étudie l'incidence,
l'homotopie ou le transport.

L'adjectif \emph{holographique} désigne ici cette compatibilité : la donnée de
frontière, accompagnée de l'orientation, de la retenue et de la mémoire, contrôle le
prolongement de l'état sans remplacer l'intérieur par un chiffre. L'unité est
la même à travers les dimensions, bien que sa décomposition et ses lecteurs
changent.

\subsection*{Holonomie nonadique et retour de phase avec accroissement de mémoire}

Le transport nonadique est une seule connexion sur le cycle des phases. Son
holonomie complète s'écrit
\begin{equation}\label{p12-83-c0-s1:eq:horizonte-conexion-nonadica}
 \Gamma_9=\operatorname{Hol}_{C_9}(\nabla_{\rm TPK}),
 \qquad g(k+9)=g(k).
\end{equation}
Ici, \(g:\mathbb Z\to C_9\), et pour chaque point
\(\widehat x_k\) d'une orbite admissible de l'état enrichi,
\[
 \Gamma_9\widehat x_k=\widehat x_{k+9},\qquad
 q_{\rm mem}(\widehat x_{k+9})
 =q_{\rm mem}(\widehat x_k)+1.
\]
En particulier, lorsque \(q_{\rm mem}\) distingue les états,
\(\widehat x_{k+9}\neq\widehat x_k\).
Les neuf composantes ordonnées décrivent un tour de la connexion ; ce ne sont pas
neuf filtres appliqués à un état immobile. Dans chaque canal, il existe une application
locale typée
\[
 \operatorname{loc}_{K,\chi}:X_\chi\longrightarrow\mathbb F_3^6,
 \qquad |\mathbb F_3^6|=729.
\]
Le transfert parcourt la fibre ambiante et localise la coordonnée de
l'extension compatible. Le tour ramène la phase et fait avancer la mémoire,
le préfixe, la retenue et la frontière.

Dans une réalisation isométrique, la séparation entre publication visible et
mémoire exprimée est typée par
\[
 J:\mathcal H_{\rm vis}\hookrightarrow\mathcal H_{\rm full},
 \quad J^*J=I,\qquad
 U_{\Gamma_9}^*U_{\Gamma_9}=I,
\]
\[
 T=J^*U_{\Gamma_9}J,\qquad
 \eta=(I-JJ^*)U_{\Gamma_9}J.
\]
On en déduit
\begin{equation}\label{p12-83-c0-s1:eq:horizonte-compresion-expresion}
 U_{\Gamma_9}J=JT+\eta,
 \qquad J^*\eta=0,
 \qquad I-T^*T=\eta^*\eta.
\end{equation}
Le terme \(\eta\) exprime exactement le défaut orthogonal produit par
cette compression. Pour des chemins composables
\(\gamma_1:x_0\to x_1\) et \(\gamma_2:x_1\to x_2\), la retenue
\(\operatorname{Carr}_R(\gamma)\in\mathsf C\) satisfait, avec
\(H(\gamma),Y(\gamma)\in\operatorname{End}(\mathsf C)\), la loi cocyclique
\begin{equation}\label{p12-83-c0-s1:eq:horizonte-carry-cociclico}
 \operatorname{Carr}_R(\gamma_2\gamma_1)
 =\operatorname{Carr}_R(\gamma_2)H(\gamma_1)
  +Y(\gamma_2)\operatorname{Carr}_R(\gamma_1),
\end{equation}
qui est la forme algébrique de la continuité généalogique.

\subsection*{Relèvement dimensionnel et coordination des régimes géométriques}

La droite marquée fournit l'ordre et l'horizon. APP construit ensuite
la carte bidimensionnelle dans laquelle somme et produit agissent sur un support
commun sans perdre leurs quotients. TRIT oriente localement chaque interaction par
la décomposition résiduelle unique
\[
 a+b=r+3c,
 \qquad a,b\in\mathbb Z,\quad
 r\in\{-1,0,+1\},\quad c\in\mathbb Z.
\]
Le Topological Phase Kernel organise chemins, feuillets, frontière, signature,
orientation, cylindre, registre, mémoire et transport. Le relèvement au cube
n'ajoute pas une métaphore spatiale : il introduit les incidences et les cohérences qui
permettent de comparer intérieur, couronne et publication dans des dimensions successives.

Cette séquence prépare également l'articulation entre trois géométries. La
structure circulaire conserve la phase ; la structure complexe représente le
quart de tour dans un espace réel \(V\) par un opérateur
\(\mathcal J:V\to V\) avec \(\mathcal J^2=-I_V\) ; la structure hyperbolique
se réalise dans le disque de Poincaré, dont la frontière est à distance
hyperbolique infinie. Leurs identifications se feront par des applications
explicites, non par égalité de dessins. L'horizon géométrique anticipe ainsi
la thèse du livre : une métrique publiée, la mémoire du parcours et
l'ouverture des continuations sont des lectures coordonnées d'une structure
orientée.

La double publication part d'une même source :
\[
 \operatorname{ev}^{\rm arq}_\infty:X_\chi\longrightarrow\mathbb R,
 \qquad
 \operatorname{ev}^{\rm inc}_\infty:X_\chi\longrightarrow\mathcal E_{\rm inc},
\]
où \(\mathcal E_{\rm inc}\) désigne la limite des objets d'incidence
de frontière construite dans la Partie II. La première application publie la
valeur archimédienne ; la seconde conserve la généalogie exceptionnelle.

\subsection*{Séquence causale APP–TRIT–TPK, continu et publications typées}

Les chapitres suivants construisent APP, TRIT et le Topological Phase Kernel
à partir de leurs données finies ; développent la structure discrète du continu et
ses deux publications ; déduisent les généalogies des constantes ; intègrent
les cinq réalisations inséparables —spectrale, solénoïdale, de jauge,
cohomologique et déterminantielle— du continu conjointement construit ; et construisent ensuite,
par la mécanique dimensionnelle, les réalisations physiques des secteurs qui
les possèdent. La publication
archimédienne apparaît à la fin de chaque chaîne, comme lecteur d'un objet déjà
construit. Tel est le sens précis dans lequel l'œuvre commence par une droite
et s'achève dans \(\mathbb R\).

\paragraph{Localisation des démonstrations complètes.}
Cette ouverture fixe les objets et l'ordre causal. Les démonstrations complètes se
trouvent dans les chapitres consacrés au complexe discret de mesure, à la
construction gnomonique d'APP, au transport orienté TRIT, à l'architecture
opératoire du Topological Phase Kernel, à la complétion cubique et à la double
publication. Les renvois ultérieurs relient chaque construction à sa
démonstration complète ; la présente ouverture fixe le domaine commun et l'ordre
de composition.


\section{Formalisation projective de l'unité à des échelles et dans des dimensions successives}
\markboth{Introduction générale}{Formalisation projective de l'unité}
\label{p12-83-c0-s2:chap:ampliacion-conservacion-unidad}

\subsection{Intervalle orienté, marques, cellules et arête de clôture}

La construction commence par une droite orientée et une échelle déclarée. On ne
commence pas par l'ensemble des nombres réels déjà complété, mais par
l'objet fini
\begin{equation}\label{p12-83-c0-s2:eq:amp-recta-k}
 \mathcal I_m=\bigl([0,10^m],\,\mathcal M_m,\,\mathcal E_m,\,o_m\bigr),
 \qquad m\geq1,
\end{equation}
où \(\mathcal M_m\) est l'ensemble des marques, \(\mathcal E_m\) celui des
intervalles élémentaires et \(o_m\) une orientation. La formule réunit quatre
objets distincts. Le segment fournit le support ordonné ; les marques
publient les positions ; les intervalles comparent les séparations ; l'orientation
décide du sens de lecture. Supprimer l'une quelconque de ces coordonnées
change le type de l'objet.

Dans la première carte, \([0,10]\), les positions nonadiques visibles sont
\(1,\ldots,9\). Entre elles existent huit intervalles intérieurs. Le pas qui
relie le dernier état visible au premier du tour suivant est
enregistré comme arête de clôture et de retenue. C'est pourquoi l'égalité
\begin{equation}\label{p12-83-c0-s2:eq:amp-nueve-marcas-ocho-intervalos}
 9\ \text{marques visibles}
 \quad\longleftrightarrow\quad
 8\ \text{intervalles intérieurs}+1\ \text{liaison de clôture}
\end{equation}
n'identifie pas les marques aux intervalles. Elle exprime comment le comptage et la mesure
partagent une frontière sans devenir la même opération.

La division euclidienne
\begin{equation}\label{p12-83-c0-s2:eq:amp-division-altura-fase}
 n=9q_9(n)+\rho_9(n),
 \qquad 0\leq\rho_9(n)\leq8,
\end{equation}
se lit désormais dans deux directions. Le résidu \(\rho_9\) détermine la phase
publiée ; le quotient \(q_9\) conserve la hauteur accumulée. Lorsque la phase
revient à sa valeur initiale après neuf pas, le quotient a avancé.
La trajectoire est hélicoïdale : une coordonnée de la roue a été récupérée,
mais pas l'état complet.

\paragraph{Interprétation algébrique de la décomposition résidu–quotient.}
L'égalité \eqref{p12-83-c0-s2:eq:amp-division-altura-fase} n'est pas seulement un algorithme pour
calculer des restes. Elle définit une extension : la projection sur le résidu contracte
l'histoire et le quotient conserve l'information nécessaire pour la relever.
Le falsificateur élémentaire consiste à imposer simultanément
\(\rho_9(n+9)=\rho_9(n)\) et \(q_9(n+9)=q_9(n)\). La première égalité est
correcte ; la seconde contredit la division euclidienne.

\subsection{Système projectif d'échelles et complétion cubique à profondeur arbitraire}

L'extension d'échelle s'effectue au moyen d'applications de raffinement, non par une
répétition graphique de la même règle. Pour \(m\leq n\), soient
\begin{equation}\label{p12-83-c0-s2:eq:amp-refinamiento-escalas}
 r_{n,m}:\mathcal I_n\longrightarrow\mathcal I_m,
 \qquad
 r_{m,m}=I,
 \qquad
 r_{n,m}r_{\ell,n}=r_{\ell,m}.
\end{equation}
L'application \(r_{n,m}\) conserve la marque publiée à la résolution \(10^{-m}\) et
contracte le détail situé en dessous de ce seuil. L'équation de
composition garantit qu'oublier deux niveaux d'un seul coup produit le même
résultat que les oublier successivement. C'est la condition qui permet
de parler d'une unique histoire observée à des échelles différentes.

En dimension trois, le passage de la fibre nonadique à la carte décimale complète
s'exprime par
\begin{equation}\label{p12-83-c0-s2:eq:amp-729-271-1000}
 1000=(9+1)^3=9^3+3\,9^2+3\,9+1
 =729+243+27+1=729+271.
\end{equation}
L'égalité contient une stratification. Les \(729\) états intérieurs
n'ont pas de coordonnée de frontière ; les \(243\) états suivants en possèdent une ;
les \(27\), deux ; le sommet en possède trois. Le nombre \(271\) compte les sommets
ajoutés par la couronne décimale. Il ne compte pas la frontière interne du cube
nonadique, dont le cardinal est \(9^3-7^3=386\), ni le complexe relatif complet,
qui inclut arêtes, faces et cubes.

Pour toute dimension \(d\), soit
\begin{equation}\label{p12-83-c0-s2:eq:amp-completacion-d}
 \overline C_9^{(d)}=(E_9\sqcup\{\star\})^d,
 \qquad \star\notin E_9,
\end{equation}
et notons \(B_r(d)\) la strate ayant exactement \(r\) coordonnées égales à
\(\star\). Alors
\begin{equation}\label{p12-83-c0-s2:eq:amp-estratos-binomial}
 |B_r(d)|=\binom dr9^{d-r},
 \qquad
 \sum_{r=0}^{d}|B_r(d)|=10^d.
\end{equation}
L'équation décrit la croissance de l'échelle sans imposer de dimension
terminale. Chaque \(d\) possède davantage d'états, mais tous appartiennent à une même loi
de complétion.

\subsection{Mesure produit compatible et conservation de l'unité sous projection}

L'expansion des cardinaux s'accompagne d'une conservation de la mesure. Attribuons
à chaque symbole de \(E_9\sqcup\{\star\}\) le poids \(1/10\) et soit \(\mu_d\) la
mesure produit sur \(\overline C_9^{(d)}\). Si
\(\pi_d:\overline C_9^{(d)}\to\overline C_9^{(d-1)}\) oublie la dernière
coordonnée, alors
\begin{equation}\label{p12-83-c0-s2:eq:amp-conservacion-proyectiva}
 (\pi_d)_*\mu_d=\mu_{d-1},
 \qquad
 \mu_d\bigl(\overline C_9^{(d)}\bigr)=1.
\end{equation}
La première égalité est la conservation évolutive de l'unité. La dimension
augmente, le nombre d'états change et la masse normalisée reste identique.
La frontière n'ajoute pas une seconde unité : elle redistribue la même unité entre
des strates compatibles.

Cette conservation possède une version cellulaire. Pour le chemin \(P_m\) et le
produit cubique \(Q_{m,d}=P_m^{\square d}\),
\begin{equation}\label{p12-83-c0-s2:eq:amp-celdas-cubicales}
 \#C_k(Q_{m,d})=\binom dk(m-1)^k m^{d-k}.
\end{equation}
En particulier, la couronne relative de \(Q_{9,3}\subset Q_{10,3}\) a pour
vecteur cellulaire
\begin{equation}\label{p12-83-c0-s2:eq:amp-corona-completa}
 \#C_\bullet(Q_{10,3},Q_{9,3})=(271,756,702,217),
 \qquad 271-756+702-217=0.
\end{equation}
La caractéristique d'Euler nulle certifie que les nouveaux sommets
s'accompagnent des incidences requises. Réduire la couronne à \(271\) élimine
précisément le transport que la théorie doit conserver.

\subsection{Courbure géométrique et fibres d'observation : séparation formelle entre métrique et horizon}

Le cinquième postulat d'Euclide régit le comportement global des
parallèles. Dans la formulation euclidienne, il fixe une géométrie dans laquelle la somme des
angles d'un triangle et la distance entre les géodésiques correspondent à une
courbure nulle. Sa modification ouvre les géométries elliptique et hyperbolique.
L'horizon visible d'une carte géométrique ne constitue pas en lui-même une
preuve du postulat : c'est la frontière où la carte cesse de distinguer les
prolongements.

La relation correcte entre les deux concepts s'établit au moyen d'applications. Une
géométrie \((M,g)\) fournit des distances et des géodésiques ; une carte
d'observation
\begin{equation}\label{p12-83-c0-s2:eq:amp-carta-horizonte}
 \operatorname{Obs}_{K}:(M,g)\longrightarrow Y_K
\end{equation}
publie l'information accessible à la résolution \(K\). L'horizon relatif
de \(x\in M\) est la fibre
\begin{equation}\label{p12-83-c0-s2:eq:amp-fibra-horizonte}
 \mathsf{Hor}_K(x)=\operatorname{Obs}_K^{-1}
 \bigl(\operatorname{Obs}_K(x)\bigr).
\end{equation}
La courbure appartient à \(g\) ; l'indiscernabilité appartient à la fibre
de \(\operatorname{Obs}_K\). L'horizon peut changer lorsque la carte est raffinée
sans que la courbure change. Cette séparation évite d'identifier une
limitation de l'observateur à une propriété intrinsèque de l'espace.

Dans le disque de Poincaré, la frontière euclidienne est à distance hyperbolique
infinie. Sur la sphère, la courbure positive ferme les géodésiques sans
frontière analogue. Dans la carte euclidienne, le parallélisme reste global. HMT
coordonne ces régimes au moyen du paramètre tritique et des algèbres
\begin{equation}\label{p12-83-c0-s2:eq:amp-tres-regimenes}
 \mathbb A_s=\mathbb R[J]/(J^2+sI),
 \qquad s\in\{+1,0,-1\},
\end{equation}
mais conserve leurs métriques et leurs codomaines. L'unité algébrique commune n'efface pas
la différence géométrique entre \(\cos\), la branche parabolique et \(\cosh\).

\subsection{Raffinement, survie et seuils de prolongement compatible}

Une expansion indéfinie contient deux orientations d'une même opération
d'échelle. Vers la précision, le raffinement distingue des cylindres de plus en
plus étroits ; vers l'extension, la hauteur enregistre des tours de plus en plus
longs. Soit \(\mathcal S_K\) l'ensemble des états survivants au niveau
\(K\). La section complète est
\begin{equation}\label{p12-83-c0-s2:eq:amp-seccion-infinita}
 X_\chi=\varprojlim_K\mathcal S_K,
 \qquad
 x=(x_K)_K,
 \qquad \tau_{K+1,K}(x_{K+1})=x_K.
\end{equation}
Chaque \(x_K\) est fini ; la compatibilité n'a pas de profondeur maximale.
Un infinitésimal ne s'identifie pas à une branche de prolongement : le premier
est une relation d'échelle dans une réalisation ; la seconde est une donnée du
système inverse. Seul un lecteur qui déclare la métrique et l'évaluation peut
transformer la seconde en le premier.

La mathématique des seuils apparaît lorsqu'un prolongement change de régime
en franchissant une frontière. Si
\(\operatorname{Surv}_K^{(N)}\) désigne les états de niveau \(K\) qui admettent une
extension jusqu'à \(N\), alors
\begin{equation}\label{p12-83-c0-s2:eq:amp-horizontes-supervivencia}
 \operatorname{Surv}_K^{(N+1)}\subseteq
 \operatorname{Surv}_K^{(N)},
 \qquad
 \operatorname{Surv}_K=\bigcap_{N\geq K}
 \operatorname{Surv}_K^{(N)}.
\end{equation}
La vacance enregistre le seuil auquel une continuation visible cesse
d'être compatible. Ce n'est pas un vide sans structure : elle conserve la cause du
rejet, la retenue et la frontière qui déterminent le niveau suivant.

\subsection{Évaluation archimédienne des sections compatibles et reconstruction de \(\mathbb R\)}

L'évaluation archimédienne apparaît après la construction de la section. Pour
une carte décimale de base \(10\), soit
\begin{equation}\label{p12-83-c0-s2:eq:amp-evaluacion-arquimediana}
 \operatorname{ev}_{\mathbb R}(x)
 =a_0+\sum_{j\geq1}a_j10^{-j},
 \qquad a_j=\operatorname{Pub}_j(x_j).
\end{equation}
La convergence s'obtient parce que les cylindres compatibles possèdent des diamètres
qui tendent vers zéro dans la carte déclarée. Le chiffre \(a_j\) est une coordonnée ;
la suite \((x_j)_j\) est l'histoire ; le réel est la publication limite.
Par conséquent,
\begin{equation}\label{p12-83-c0-s2:eq:amp-cadena-recta-R}
 [0,10]\longrightarrow[0,1000]
 \longrightarrow[0,10^m]
 \longrightarrow X_\chi
 \xrightarrow{\operatorname{ev}_{\mathbb R}}\mathbb R
\end{equation}
est une chaîne de types, non une chaîne d'inclusions naïves.

La seconde publication du même \(x\) conserve l'incidence exceptionnelle.
L'œuvre parcourra les deux directions et reviendra finalement à
\(\mathbb R\). Le retour sera alors une conclusion : la droite réelle aura
été expliquée comme une évaluation de l'état compatible, et non utilisée
comme présupposé rendant la généalogie inutile.


\section{Généalogie historique du problème du continu et formulation mathématique HMT–MD}
\markboth{Généalogie historique du continu}{Généalogie historique du continu}

\subsection*{Le problème historique du continu, de la valeur et de la mesure}

Le continu constitue l'un des problèmes fondateurs des
mathématiques. L'arithmétique représente des unités discrètes ; la géométrie
compare des étendues ; l'analyse complète des suites et construit des limites ; la
théorie de la mesure attribue des grandeurs à des ensembles ; la physique transforme ces
grandeurs en observables. L'efficacité de ces disciplines repose sur la
possibilité de publier un résultat stable et reproductible. Le problème qui
gouverne cette œuvre apparaît un niveau avant : déterminer quelle structure rend
possible cette stabilité et quelle information doit être conservée pour récupérer une
valeur, une mesure et leurs symétries à partir d'une origine commune.

Sur la droite discrète, une marque et l'espace qui la sépare de la suivante sont des
objets distincts. La marque répond à l'opération de compter ; l'intervalle,
à l'opération de mesurer. Le passage de \(0\) à \(10\) contient les neuf marques
visibles \(1,\ldots,9\), les huit intervalles intérieurs entre marques
consécutives et une arête relevée de clôture qui transporte le comptage au
tour suivant. La lecture modulaire conserve la phase ; le quotient conserve
la hauteur atteinte. Lorsque les deux coordonnées sont maintenues ensemble, la
répétition circulaire devient une hélice et la suite des entiers
conserve l'histoire de ses retours.

L'holographie modulaire triadique (HMT) développe cette intuition comme une
construction mathématique. Un système fini de positions, de feuillets,
d'orientations, de chemins et de transports produit des familles compatibles d'états.
Ses évaluations publient des chiffres et des valeurs ; ses relèvements conservent
résidu, quotient, retenue, frontière et mémoire. Le continu se réalise comme
l'image d'une section compatible du système des prolongements. Sa
structure discrète réside dans la généalogie qui relie toutes les
résolutions.

Cette formulation distingue trois niveaux. Un \emph{chiffre} est une coordonnée
visible produite par un lecteur. Une \emph{valeur} est l'image d'une section
sous une carte donnée. Un \emph{nombre généalogique} est la section
complète qui conserve profondeur, opération, orientation, feuillet, résidu,
quotient, retenue, chemin, frontière, mémoire et survie. Le développement
décimal constitue une évaluation de cet objet ; la compatibilité entre ses
étapes constitue sa continuité.

La même structure transforme le concept de symétrie. Une symétrie est la
conservation évolutive de l'identité sous une transformation avec mémoire.
Un cycle peut retrouver la position visible et déplacer simultanément
l'état vers une autre hauteur, un autre feuillet ou un nouveau bloc. L'invariance appartient à
la généalogie complète. La valeur réelle et la symétrie exceptionnelle apparaissent,
par conséquent, comme deux projections coordonnées d'un même antécédent :
l'une contracte des cylindres compatibles et l'autre conserve l'incidence que la
contraction archimédienne laisse hors de la coordonnée publiée.

La mécanique dimensionnelle (MD) prolonge cette construction dans le domaine physique.
La section numérique reçoit une représentation spectrale ; la persistance se
réalise comme particule-chemin ; le registre du chemin produit une signature ; la
signature alimente un caractère et une hiérarchie de tours ; la carte dimensionnelle
publie finalement une grandeur. La chaîne directrice de l'œuvre est
\[
 \text{unité}\longrightarrow\text{état avec mémoire}
 \longrightarrow\text{continu, nombre et symétrie}
 \longrightarrow\text{particule, constante et grandeur}.
\]

L'objectif scientifique est unitaire et se développe en cinq Parties. Le
noyau formel construit APP, TRIT et le Topological Phase Kernel ; la
structure discrète du continu déploie ses sections, la double
projection et la généalogie exceptionnelle ; la généalogie des constantes
produit \(\pi\), \(e\), \(\varphi\), \(\alpha\), l'échelle interne
d'action et ses transducteurs ; cinq chaînes de spécialisation construisent leurs
antécédents internes et isolent les comparateurs qui conduisent à chaque
reconnaissance classique ; la mécanique dimensionnelle construit ensuite les
constantes constitutives, la masse, la statistique quantique, l'information et
la gravité. Chaque mesure est présentée avec l'objet et l'opération
qui l'engendrent ; chaque confrontation métrologique s'ouvre après la fixation de la
sortie interne.

L'ordre de l'œuvre suit la dépendance constructive : un objet apparaît là
où il commence à gouverner les résultats ultérieurs. Chaque énoncé déclare
ses prémisses, sa démonstration et, le cas échéant, le morphisme de
réalisation physique. De cette manière, une comparaison métrologique ultérieure ne
remplace pas la généalogie qui a produit la sortie.

\section*{De la clôture généalogique à la réalisation physique}

La thèse de cet ouvrage acquiert quatre clôtures coordonnées et irréductibles :
généalogique, opératorielle, matricielle et physique. Il ne s'agit pas de juxtaposer une
théorie des ensembles, une théorie des algèbres d'opérateurs ou un formalisme
quantique à une construction numérique préalable. Les quatre développements naissent de
la même préséance causale : l'APP construit la double réalisation positionnelle et
arithmétique ; le TRIT orienté distingue retour, présent torsionnel et ouverture ;
le TPK conserve valeur, feuillet, orientation, résidu, quotient, retenue, bord,
route, registre, mémoire, survie, phase et jauge. Ce n'est qu'ensuite que se forment
la section compatible, ses lecteurs et ses représentations. Le diagramme directeur
est donc
\[
\begin{gathered}
 \mathrm{APP}
 \longrightarrow \mathrm{TRIT}_{\rm orientado}
 \longrightarrow \mathrm{TPK}
 \longrightarrow \mathcal X_\infty^{\rm cal},\\[2mm]
 \mathcal X_\infty^{\rm cal}
 \xrightarrow{\;\mathfrak R\;}\mathbb R,
 \qquad
 \mathcal X_\infty^{\rm cal}\times\mathcal C_{\rm exc}
 \xrightarrow{\;\mathfrak E\;}\mathcal E_{\rm exc},\\[2mm]
 C(X_m)\xrightarrow{\;\operatorname{diag}\;}M_{9^m}(\mathbb C)
 \xrightarrow{\;a\mapsto a\otimes I_9\;}
 \mathrm{UHF}(9^\infty)
 \xrightarrow{\;\pi_\tau(\,\cdot\,)''\;}R_{\rm hyp},\\[2mm]
 \mathsf P_{\rm CT}^{+}
 \xrightarrow{\;\mathfrak R_{\rm MD}^{\rm disc,+}\;}
 \mathbf{SpecRoute}\times\mathsf{Hilb}_{\mathbb R,J}
 \times\mathbf B\mathfrak M_D^+\times\mathbf B G_C.
\end{gathered}
\]
Les quatre lignes ne forment pas une unique chaîne linéaire. L'évaluation réelle et la
projection exceptionnelle sont des lectrices sœurs de la section enrichie ; la
branche traciale passe par les algèbres cylindriques et leurs inclusions unitaires ; la
réalisation physique part de routes préparées munies des structures
exigées par chaque composante. Cette séparation des domaines empêche de déduire le facteur de
von Neumann directement d'une section, d'identifier la projection
exceptionnelle à une action sur les chiffres ou de présenter une grandeur physique comme
s'il s'agissait d'un état TPK sans morphisme de réalisation.

\subsection*{Clôture généalogique : extensionnalité, infini et décision}

La première clôture détermine ce que signifie conserver un objet à travers
toutes les profondeurs. Étant donné un système inverse d'espaces finis
\((X_m,p_m)\), une section \(x=(x_m)_m\) n'est pas caractérisée uniquement par
l'égalité de ses valeurs publiées. Le théorème de descente par fibres fixe
les conditions dans lesquelles un lecteur fini descend à un quotient et
sépare surjectivité, constance sur les fibres et unicité de l'application descendue.
Le lemme de König garantit l'existence d'une branche lorsque les niveaux sont
finis, non vides et finement ramifiés ; l'unicité exige, en outre, le
caractère complet qui sélectionne un unique prolongement compatible.

Sur une section fixée, les cylindres décroissants fournissent une
réalisation exacte du passage \(0\!-1\!-\!\infty\) : le support se contracte, la
hauteur augmente, la masse demeure unitaire et les mesures cylindriques convergent
faiblement-* vers \(\delta_x\). Ce résultat coordonne le delta de Kronecker fini,
les espérances conditionnelles entre niveaux et le delta de Dirac comme mesure
ponctuelle, sans identifier ce dernier à un opérateur différentiel. La
fondation des préfixes empêche les chaînes descendantes infinies dans chaque
présentation finie ; la coinduction contrôle le prolongement compatible. Les
ordinaux de von Neumann sont relevés avec une fibre de mémoire de route, et la
cardinalité apparaît comme une décatégorification qui ne retient pas toute la généalogie.

La même distinction oblige à séparer puissance cylindrique, puissance fermée et
puissance pleine, ainsi que choix ensembliste, choix naturel, choix
algébrique et choix effectif. Le codage dans ZF/ZFC conserve les
constructions et permet de les confronter à l'extensionnalité, à Fondation,
à Puissance et à Choix ; il ne fait pas de la droite réelle ni de l'univers ensembliste
l'origine d'APP--TRIT--TPK\@. Dans le domaine effectif, la décidabilité de chaque
fenêtre finie, l'existence de la limite, l'unicité d'une instance et
l'impossibilité d'un décideur total uniforme sont quatre assertions distinctes.
La réduction de Godel--Turing n'affecte que la quatrième : elle n'invalide pas les sélecteurs
particuliers déjà construits. De même, la confrontation Godel--Cohen sépare
l'extension \(M\mapsto M[G]\) de l'enrichissement interne d'un état et
évite de confondre le franchissement des denses de profondeur avec la généricité sur un
modèle. Les cylindres boréliens, les hiérarchies de comparaison et le passage
projectif de minimax sont introduits avec leurs hypothèses de compacité,
de continuité et de cohérence, non comme des conséquences automatiques de la notation
inverse.

Cette clôture possède également une formulation informationnelle. Si une mise à jour
\(U\) de l'état complet est bijective, alors
\(H(X\mid U(X))=0\), et le terme logique idéal de Landauer s'annule. Si un
lecteur \(q\) ne conserve que l'ombre visible, l'information écartée est
localisée dans \(H(X\mid q(X))\). Dans la formulation opératorielle, un automorphisme
préservant la trace conserve l'entropie, tandis qu'une espérance conditionnelle
préservant la trace \(E_N:M\to N\) satisfait
\[
 S_\tau(E_Nh)-S_\tau(h)=D_\tau(h\Vert E_Nh)\geq0.
\]
Le coût correspond à l'oubli de la fibre, non au transport réversible de
l'état enrichi ; cette identité n'affirme ni une dissipation physique totale nulle ni
une complexité computationnelle réduite.

\subsection*{Clôture opératorielle : de l'unité nonadique à la dimension continue}

La seconde clôture réalise opératoriellement la conservation évolutive de
l'unité. Pour le système complet, on considère
\[
 A_m=M_{9^m}(\mathbb C),
 \qquad \iota_m(a)=a\otimes I_9,
 \qquad
 \tau_m(a)=9^{-m}\operatorname{Tr}(a).
\]
L'identité \(\tau_{m+1}\circ\iota_m=\tau_m\) conserve simultanément
l'unité et la trace. La limite inductive est l'UHF de type surnaturel
\(9^\infty=3^\infty\), et la clôture faible de sa représentation GNS traciale est
le facteur hyperfini \(R_{\rm hyp}\) de type \(II_1\). Les valeurs
\(\operatorname{rank}(P)/9^m\) se densifient et la trace des projecteurs dans
\(R_{\rm hyp}\) réalise tout l'intervalle \([0,1]\) : une dimension continue
émerge d'une tour d'inclusions discrètes qui préservent l'unité.

Cette construction ne se confond pas avec la branche marquée
\(j_{m,r}(a)=a\otimes p_r\). Les isométries d'une fibre individuelle engendrent en
norme les compacts sur la limite hilbertienne et, en clôture faible, un facteur
de type \(I_\infty\). La sélection d'un secteur conserve la marque, mais non
l'unité de la tour ; c'est le relèvement simultané des neuf fibres qui
produit la branche UHF--\(II_1\). On n'identifie pas non plus sans preuve la réalisation
UHF à la réalisation par produit croisé de l'horloge profinie : une
conjugaison canonique exige de déclarer l'action, la mesure, la moyennabilité,
l'ergodicité et l'isotropie. La logique orthomodulaire des projecteurs, les contextes
booléens, les états normaux et la normalisation entropique
\[
 S_\tau(9^m\rho)=S_{\rm vN}(\rho)-m\log9
\]
sont dérivés au sein de cette construction opératorielle et maintiennent séparés
l'état matriciel, sa densité traciale et sa publication statistique.

La topologie, la dimension et la composition ne se confondent pas dans le facteur. La clôture
transversale est enregistrée au moyen de
\[
 \mathfrak C_{\rm HMT}
 =\bigl(C(\Omega_\infty),R_{\rm hyp},\mathcal G_{\rm TPK}\bigr).
\]
Le spectre commutatif conserve les points et les cylindres ; le facteur tracial conserve
la dimension et la statistique ; le groupoïde conserve la concaténation, l'inversion
admissible et la provenance des routes. Cette architecture triple est formulée après
la double projection et le corridor exceptionnel. Elle ne les fonde ni ne les
remplace : l'évaluation archimédienne et la chaîne
Paley--Witt--Golay--Leech--\(V^\natural\)--Monstre--Moonshine demeurent
les deux lectures coordonnées du même antécédent enrichi.

\subsection*{Clôture matricielle : profondeur, taille extérieure et dilatation}

La troisième clôture introduit un contrôle simultané de trois indices. Si
\(m\) mesure la profondeur généalogique, \(n\) la taille extérieure d'une grille
et \(s\) le niveau matriciel de ses entrées, la famille
\(\mathfrak M_{m,s}^{(n)}\) est bigraduée par \((m,s)\) lorsque \(n\) demeure
fixe et trigraduée par \((m,n,s)\) lorsque la taille extérieure varie aussi.
En particulier,
\[
 \mathcal K_{m,s}
 :=\operatorname{UCP}\!\bigl(C(X_m),M_s(\mathbb C)\bigr)
\]
possède deux opérations indépendantes : la restriction généalogique
\(\rho_m^s(\Phi)=\Phi\circ\jmath_m\) et la compression matricielle
\[
 \operatorname{mc}_{\boldsymbol V}(\Phi_1,\ldots,\Phi_k)(f)
 =\sum_rV_r^*\Phi_r(f)V_r,
 \qquad \sum_rV_r^*V_r=I_s.
\]
Les deux commutent exactement :
\[
 \rho_m^s\!\left(\operatorname{mc}_{\boldsymbol V}(\Phi_1,\ldots,\Phi_k)\right)
 =\operatorname{mc}_{\boldsymbol V}\!\left(
   \rho_m^{s_1}(\Phi_1),\ldots,\rho_m^{s_k}(\Phi_k)\right).
\]
L'identité est fonctorielle et distingue raffinement de l'histoire et réduction
de l'observateur.

Toute famille compatible en profondeur détermine une unique application UCP
sur \(C(X_\infty)\) et admet une dilatation de Stinespring commune. Dans le
corridor diagonal, une unique mesure spectrale projective \(P\) satisfait
\[
 E_x^{(m)}=V^*P([x]_m)V
\]
pour toutes les profondeurs et tous les cylindres. La dilatation minimale est
unique à équivalence unitaire près et conserve exactement ce qui est vu par le
lecteur ; une somme directe permet une réalisation environnementale fidèle qui conserve
toutes les distinctions de l'algèbre. Sur la limite UHF, une famille compatible
d'applications UCP possède également un prolongement et une dilatation communs.
L'espérance conditionnelle
\[
 \mathbb E_m=\operatorname{id}\otimes\tau_9:
 M_{9^{m+1}}(\mathbb C)\longrightarrow M_{9^m}(\mathbb C)
\]
préserve la trace et écarte le dernier facteur de manière contrôlée ; son environnement
de Stinespring enregistre le chiffre omis. En général, il n'existe pas de mémoire
auxiliaire finie uniforme pour toutes les profondeurs. De plus, la fidélité d'une
représentation n'implique pas à elle seule la fidélité dynamique APP--TPK : celle-ci
exige un opérateur d'entrelacement compatible avec la route, la monodromie et la composition.

La théorie matricielle se concrétise en objets vérifiables. Les carrés magiques
quantiques qutrit de paramètres \((9,3)\) et \((12,3)\) possèdent respectivement
81 et 144 effets, des sommes de lignes et de colonnes exactement unitaires et des
dilatations simultanées déclarées ; les mêmes POVM produisent des cubes magiques
quantiques d'ordres neuf et douze. Le cas d'ordre douze est organisé au moyen de
la fibration \(12\to4\times3\), avec sa jauge et sa condition explicite
d'équivariance. Par une autre voie, les espérances conditionnelles somme et produit de l'APP
engendrent l'algèbre
\[
 C^*(P_{\Sigma,2},P_{\Pi,2})
 \cong\mathbb C^4\oplus M_2(\mathbb C)\oplus M_2(\mathbb C),
\]
et le bloc central réalise à chaque niveau \(s\) l'intervalle matriciel libre
\([5I_s,9I_s]\). Ces constructions distinguent carrés latins quantiques,
carrés magiques, mesures projectives, mesures positives et semi-classicité ;
aucune identification n'est obtenue uniquement parce que les objets partagent l'ordre
neuf.

L'antécédent spectral commun de deux lecteurs n'est formulé qu'après
les avoir construits tous les deux. Pour des applications finies
\(R_m:X_m\to\mathcal A_m\) et \(Q_m:X_m\to\mathcal E_m\), les incidences
\[
 P_{a,e}=1_{\{R_m=a\}}1_{\{Q_m=e\}}
\]
raffinent les deux mesures spectrales diagonales et récupèrent leurs marginales. Le
résultat ne prouve pas l'injectivité conjointe et n'impose pas qu'un lecteur soit
une compression de l'autre. Dans le régime non commutatif, deux PVM qutrit mutuellement
non biaisées ne possèdent pas de PVM conjointe nette ; la réalisation simultanée exige
un bruit compatible ou un environnement plus grand. Cette différence de type préserve la
double projection HMT : le théorème matriciel classe des réalisations de lecteurs
et ne réécrit pas leur antécédent fondateur.

\subsection*{Clôture physique : domaine commun, deux orientations et localisation}

La quatrième clôture construit le domaine depuis lequel une route TPK peut recevoir
simultanément une réalisation spectrale, hilbertienne, dimensionnelle et de Cartan.
Soient \(\mathsf D_{\rm sp}\), \(\mathsf D_Q\), \(\mathsf D_D\) et
\(\mathsf D_C\) les domaines de référence de ces quatre représentations, avec des
foncteurs d'oubli fidèles vers la catégorie
\(\mathsf P_{\rm TPK}^{\rm enr}\) des états préparés et des routes enregistrées.
Le produit fibré
\[
 \mathsf P_{\rm CT}^{+}
 :=\mathsf D_{\rm sp}
 \times_{\mathsf P_{\rm TPK}^{\rm enr}}\mathsf D_Q
 \times_{\mathsf P_{\rm TPK}^{\rm enr}}\mathsf D_D
 \times_{\mathsf P_{\rm TPK}^{\rm enr}}\mathsf D_C
\]
impose que les quatre composantes aient exactement le même antécédent.
Ses projections canoniques sont des adaptateurs fonctoriels et produisent
\[
 \mathfrak R_{\rm MD}^{\rm disc,+}
 =\bigl(\mathfrak S_+,\mathcal Q_+,\mathfrak D_+,
        \rho_{C,+}^{\rm disc}\bigr),
\]
à valeurs, respectivement, dans les routes spectrales, les espaces de Hilbert réels
à structure complexe, le monoïde dimensionnel positif et le groupe de
Poincaré discret. La construction affirme le foncteur sur le domaine commun ; elle
n'affirme pas que toute route TPK admette automatiquement les quatre relèvements.

Sur une route \(\gamma=e_1\cdots e_r\), la variation totale
\[
 V_+(\gamma)=\sum_jw(e_j)
\]
et le transport net orienté
\[
 V_{\rm or}(\gamma)=\sum_j\varepsilon(e_j)w(\underline e_j)
\]
sont des lecteurs frères. La boucle \(ee^\vee\) démontre que le premier ne se
factorise pas, en général, par le second : elle possède un transport orienté nul et une
variation positive. Pour obtenir une réalisation inversible, on restreint le
domaine aux arêtes dont les composantes spectrale et hilbertienne sont
inversibles, dont la donnée dimensionnelle est une cochaîne additive et dont les cocycles de
Cartan sont normalisés et antisymétriques. Ce n'est qu'alors que la localisation
\(\Pi_\vee(\mathsf P_{{\rm CT},{\rm rev}}^+)\) produit le foncteur orienté
\(\mathfrak R_{\rm MD}^{\rm disc,or}\) et transporte l'inversion composante par
composante. La lecture positive n'est pas éliminée et la lecture orientée n'est pas obtenue
en changeant uniquement un signe.

Cette construction physique précède les constructions de Planck et de l'électron
parce qu'elle leur fournit domaine, représentations et règles de transport, mais
n'absorbe pas leurs démonstrations de référence ni la loi générale des masses. La
chaîne nombre--particule--route--registre chronologique enrichi--signature--
caractère--tours conserve sa direction causale ; l'échelle absolue d'action,
l'électron central, les familles de masse, CKM, PMNS et les secteurs physiques
maintiennent leurs objets, domaines et conclusions propres. Une coïncidence métrologique n'est
confrontée qu'après avoir figé la sortie interne. Les extensions qui
exigent encore un opérateur d'entrelacement, une limite lisse, une preuve de stabilité dynamique
ou un morphisme physique sont énoncées avec ces données comme hypothèses exactes.
L'absence de l'une de ces applications empêche d'étendre le théorème au codomaine
correspondant, sans altérer les théorèmes généalogiques, opératoriels et
matriciels déjà démontrés dans leurs domaines.

La lecture de l'ouvrage suit, par conséquent, l'ordre de dépendance et non
celui de la reconnaissance extérieure :
\[
\begin{aligned}
 \mathrm{APP}&\longrightarrow\mathrm{TRIT}_{\rm orientado}
 \longrightarrow\mathrm{TPK}\\
 &\longrightarrow\text{section enrichie}
 \longrightarrow
 \begin{cases}
  \text{évaluation archimédienne},\\
  \text{projection exceptionnelle},\\
  \text{algèbres cylindriques et clôture traciale},\\
  \text{routes préparées et réalisation physique}.
 \end{cases}
\end{aligned}
\]
Chaque sortie conserve son domaine, ses hypothèses et ses obstructions précises. Ce qui est commun n'est pas
une égalité forcée entre codomaines, mais la provenance du même
état enrichi. Tel est le sens précis de la clôture holographique : aucune
publication terminale n'épuise la généalogie qui la produit et aucune clôture
ultérieure ne peut remplacer la double direction --archimédienne et exceptionnelle--
qui conserve l'unité avec mémoire.


\subsection*{Perspective historique et développement de la théorie}

L'histoire des mathématiques peut se lire comme un élargissement successif des
opérations admises sur l'unité. L'entier a rendu possible le comptage ;
le rationnel a incorporé la division ; le réel a clos les processus
d'approximation ; le complexe a fait de la rotation une opération algébrique.
Chaque élargissement a fixé un domaine, une équivalence et une représentation.
L'égalité a publié le résultat avec une immense efficacité et a absorbé, en même temps, une bonne
partie du chemin concret qui l'avait produit.

La géométrie a suivi un développement analogue. La construction euclidienne a fixé
l'incidence et la mesure ; les géométries non euclidiennes ont montré la coexistence
de régimes de courbure formellement cohérents
\cite{EuclidElements,PoincareScienceHypothesis}. Le cercle trigonométrique,
le plan complexe et le disque hyperbolique coordonnent rotation, périodicité et
transport au moyen d'espaces d'états différents. HMT situe sous ces
réalisations un état ternaire commun et construit à partir de lui les branches
elliptique, parabolique et hyperbolique.

La mécanique a incorporé à l'unité une dépendance relationnelle. La loi du
levier a relié poids, bras et équilibre ; la critique de Mach du mouvement
absolu a placé au premier plan la relation entre corps, référentiel et totalité
\cite{ArchimedesEquilibrium,MachMechanics} ; la relativité a formalisé les
invariances cinématiques, l'équivalence locale et la géométrie du champ
gravitationnel~\cite{Einstein1918Principles} ; la théorie quantique a introduit une
échelle d'action et une propagation par amplitudes de trajectoire
\cite{PlanckNobelLecture,Feynman1948}. Le programme HMT--MD réunit ces
exigences dans un problème préalable : construire l'état qui conserve position,
orientation, opération, mémoire et prolongement avant de publier un nombre ou
une grandeur.

Les trois termes du nom précisent la méthode. \emph{Holographie} désigne
la conservation de la généalogie et de la règle de prolongement entre
résolutions. \emph{Modulaire} désigne la coordination de la carte modulo neuf,
de l'orientation modulo trois, des lecteurs résidu--quotient et de la publication
décimale en base mille. \emph{Triadique} désigne aussi bien l'orientation locale
\(\{-1,0,+1\}\) que l'organisation des niveaux qui relient le support
fini, la section numérique et sa réalisation physique.

\subsection*{Conservation évolutive et signification de la symétrie}

HMT formule la symétrie comme conservation de l'identité sous transformation
avec mémoire. Une opération peut retrouver la même phase visible sans
retrouver le même état enrichi. Si \(\Theta\) désigne le transport,
la relation
\[
 \operatorname{pr}_{\rm vis}\circ\Theta^9
 =\operatorname{pr}_{\rm vis}
\]
coexiste avec un accroissement non trivial de retenue ou de mémoire. La répétition
horizontale est la projection d'une évolution hélicoïdale. L'identité est
conservée parce que sa généalogie permet de la reconnaître à travers le changement, non
parce qu'elle reste immobile.

L'extrême et moyenne raison holographique réalise ce principe dans la
complétion cubique
\[
 999=729+243+27,\qquad
 1000=729+243+27+1.
\]
Le sommet \(+1\) de la complétion cubique et le terme terminal \(+I\) de la
clôture exponentielle sont des opérations de clôture définies dans des domaines
distincts. Leur correspondance structurelle réside dans le fait que toutes deux achèvent une
évolution et conservent l'unité dans leur codomaine respectif. Cette conservation
organise le retour nonadique, la retenue décimale, la croissance intérieure des
cylindres et la croissance extérieure des incidences.

Le prolongement dimensionnel s'exprime par
\[
 10^d=(9+1)^d=\sum_{k=0}^{d}\binom dk9^k.
\]
Après normalisation par \(10^d\), les coefficients forment une famille de mesures
projectivement compatible dont la masse totale reste égale à un.
L'intérieur, les strates de frontière et le sommet redistribuent cette unité lors du
changement de rang. Dans la réalisation cubique
\(\mathcal C_9=(\mathbb Z/9\mathbb Z)^3\), les six projections orientées
conservent le support APP et l'opérateur face--sommet produit un espace de
rang quatre. La signature \((3,1)\) émerge de cet opérateur et fournit
l'antécédent fini de la lecture spatio-temporelle ultérieure. Dimension,
frontière et conservation de l'unité sont ici des résultats du même système
de projections.

\subsection*{Unification algébrique des géométries elliptique, parabolique et hyperbolique}

Le point de départ conceptuel était d'articuler trois réalisations que la
mathématique présente habituellement dans des espaces différenciés : le plan
complexe, le cercle trigonométrique et le disque hyperbolique de Poincaré. HMT
les réunit au moyen de la famille
\[
 \mathbb A_s=\mathbb R[X]/(X^2+s),
 \qquad s\in\{+1,0,-1\}.
\]
La branche elliptique contient une racine interne de \(-1\) et la clôture circulaire ; la
branche parabolique contient le seuil nilpotent ; la branche hyperbolique contient
les deux idempotents et l'ouverture de Poincaré. Ces trois algèbres apparaissent
dès TRIT et gouvernent ensuite les feuillets, les canaux, la métrique de
Lorentz--Minkowski et la lecture physique du temps.

La différence entre somme et produit fournit l'opération élémentaire qui
articule ces géométries. La somme accumule déplacements et alternatives ; le
produit compose des étapes et replie les entrées sur une même classe lorsque
des éléments non inversibles interviennent. APP conserve les deux évaluations sur
un support commun. Sa séparation spectrale et les projecteurs \(P_\pm\) du
bloc fixent les modes propres ; les projecteurs conditionnels
\(P_{\Sigma,d}\) et \(P_{\Pi,d}\), avec leur commutateur, construisent un précurseur
fini de l'incompatibilité de Heisenberg. La
lecture physique ultérieure transforme cette séparation en coordonnées
angulaires conjuguées \(A\) et \(C^*\),
en canaux orientés et en holonomie.

\subsection*{Construction du noyau formel}

APP se construit à partir des positions \(1,\ldots,9\) en deux coordonnées. Le
noyau \(8\times8\), la publication de
la classe nulle par \(9\) et le gnomon de dix-sept cellules complètent la
carte \(9\times9\). L'identification des côtés opposés produit le tore
\[
 X_9=(\mathbb Z/9\mathbb Z)^2
\]
et ses déplacements engendrent le graphe de Cayley. Les évaluations
\(\Sigma\) et \(\Pi\) agissent sur les mêmes sommets et chemins.

TRIT est l'état ternaire orienté. Dans un relèvement
\[
 a+b=r+3c,
\]
le résidu \(r\) publie la sortie et le quotient \(c\) conserve la manière dont elle a été
obtenue. Le triplet distingue orientation, seuil et mémoire. Le TPK est le
système opératoire de dynamique généalogique qui incorpore le sélecteur tritique
du mode opératoire, transporte l'état et conserve simultanément les deux
feuillets APP avec leur position, orientation, phase, retenue, frontière, chemin,
registre et survie. Ses lecteurs
nonadique, ternaire et positionnel, son système de douze phases et ses relèvements
forment un unique arbre de domaines et de codomaines. Les dépendances mathématiques
sont les flèches typées de cet arbre ; chaque application est présentée avec son domaine, son
codomaine et un calcul reproductible.

Plus précisément, APP est le quintuplet
\[
 \mathcal A_9=
 \bigl(X_9,\Gamma_{\rm APP},\Sigma,\Pi,
       \operatorname{Path}(\Gamma_{\rm APP})\bigr),
 \qquad
 \Gamma_{\rm APP}=\operatorname{Cay}
 \bigl(X_9,\{\pm e_1,\pm e_2\}\bigr),
\]
avec évaluations
\[
 \Sigma(i,j)=\operatorname{dr}_9(i+j),\qquad
 \Pi(i,j)=\operatorname{dr}_9(ij),
 \qquad
 \operatorname{dr}_9(n)=1+((n-1)\bmod9).
\]
Le relèvement conserve simultanément
\(i+j=R^+_{ij}+9Q^+_{ij}\) et
\(ij=R^\times_{ij}+9Q^\times_{ij}\) ; la valeur publiée et l'information
de quotient appartiennent à des coordonnées différentes d'un même état.

Le relèvement de TRIT, \(\widehat\tau=(r,c)\), conserve résidu et retenue. La fibre
au-dessus du symbole visible \(0\) contient \((0,-1)\), \((0,0)\) et
\((0,+1)\), dont les arbres de prolongement peuvent différer. Cette multiplicité
est la forme minimale sous laquelle la même coordonnée visible soutient des histoires
distinctes.

Le TPK s'organise en huit classes opératoires avant de présenter ses
réalisations particulières : support et état ; sélection interne ; transport ;
lecteurs et quotients ; mémoire et frontière ; périodicité et retour ;
prolongement ; et sorties. Cette classification sépare les objets du porteur,
les applications qui mettent à jour l'état, les lecteurs qui publient ses coordonnées
et les foncteurs qui contractent la section vers un codomaine archimédien,
exceptionnel ou classificatoire. Chaque opérateur du TPK reçoit ainsi un domaine, un
codomaine, une règle et ses invariants avant de recevoir un nom abrégé.

Les phases nonadiques se répartissent en trois fibres opératoires :
\(\{1,4,7\}\) active l'évaluation additive, \(\{2,5,8\}\) active
la multiplicative et \(\{3,6,9\}\) enregistre le seuil. La transition
\(9\to1\) restitue la phase publiée et élève la mémoire dodécaphasique ; le
cycle visible est donc la projection d'une hélice d'états
enrichis.

Une fois les trois objets présentés, l'œuvre développe leurs conséquences.
Le bloc central
\[
 B_c=9P_++5P_-
\]
fixe le coefficient local d'échange \(2\) et l'écart spectral
\(9-5=4\) ; l'agrégation modulaire produit \(51-45=6\), et le déploiement sur
les neuf positions produit \(54=9\cdot6\). La forme de signature mixte et
\(SO^+(1,1)\) construisent la géométrie lorentzienne finie. La complétion
cubique construit l'extrême et moyenne raison holographique, la loi d'aire
discrète et la frontière de propagation.

\subsection*{Algorithme uniforme en profondeur finie et publication coinductive
corrélée de \texorpdfstring{\(\pi,e,\varphi,\alpha\)}{pi, e, phi, alpha}}

La génération numérique est régie par l'entonnoir
\[
 104\,976\to468\,U_6\to243\,w_6
 \to w_{30}\to R_{36}\to G_9.
\]
Les canaux de \(\pi\), \(e\), \(\varphi\) et \(\alpha\) conservent des lecteurs
typés différents : clôture angulaire, propagation autoproportionnelle,
autoéchelle et couplage signé dodécaphasique. La relation interne
\(\operatorname{Ext}\), la survie coinductive et la partition
exhaustive de la fibre produisent à chaque niveau un unique cylindre prolongeable
et compatible avec la troncature précédente ; la famille emboîtée détermine
une section enrichie unique de toute profondeur. Ce n'est qu'ensuite que les
encadrements de Leibniz, le semi-groupe factoriel et Perron--Fibonacci localisent
archimédiennement les blocs déjà engendrés. Les exécutions de mille et de dix mille
chiffres sont des certificats finis étendus de cette construction uniforme, non son
fondement logique.

Le prolongement initial construit
\[
 w_6\longrightarrow w_{12}\longrightarrow w_{18}
 \longrightarrow w_{24}\longrightarrow w_{30}
 \longrightarrow R_{36},
\]
et le recensement interne sélectionne avant l'évaluation
\[
 w_\pi=010211,\qquad w_e^{(\mathrm{Euler})}=201101,\qquad
 w_\varphi=121200.
\]
L'exposant précise que \(w_e^{(\mathrm{Euler})}\) appartient au lecteur du
nombre d'Euler \(e\) ; c'est un objet distinct du chemin électronique
\(\Gamma_e\) de la mécanique dimensionnelle.
Pour chaque profondeur finie prescrite, \(\operatorname{Ext}\), le
transport TPK et la survie compatible produisent d'abord exactement un
cylindre \(C_{\chi,n}\) parmi les \(729\) raffinements. Le lecteur rationnel
s'applique ensuite au cylindre déjà engendré : il construit son encadrement interne et
vérifie \(j_{\min}=j_{\max}\) comme certificat archimédien de localisation,
sans sélectionner la sortie ni fournir de chiffres cibles au générateur. Le
\cref{thm:generacion-monodromica-conjunta-rev7} démontre, pour
\(\chi\in\{\pi,e,\varphi\}\), la non-vacuité, l'unicité et la
compatibilité de ces survivants à toute profondeur finie ; par conséquent,
les partitions construisent
\[
 c_\chi=(C_{\chi,n})_{n\ge0}
 \in\varprojlim_n\{C_{\chi,n}\}.
\]
Les cylindres sont emboîtés, leurs diamètres convergent vers zéro et le
relèvement EF--\(G_9\) conserve la fibre enrichie. Ainsi sont
séparés l'algorithme uniforme paramétré par la profondeur, la section
coinductive qu'il détermine et ses exécutions finies de mille et de plus de dix
mille chiffres, qui agissent comme certificats étendus et non comme fondement
logique du théorème.

Les trois sections remplissent des fonctions structurelles différenciées. \(\pi\)
organise clôture, orientation et frontière monodromique ; \(e\) organise la
propagation autoproportionnelle et l'évolution du préfixe ; \(\varphi\) organise
l'autoéchelle de Perron--Fibonacci et son mode stable \(\varphi^{-2}\).
L'opérateur orienté \(\mathscr R(x,y)=x/y^2\) publie les lectures conjuguées
\(\pi/\varphi^2\) et \(\varphi/\pi^2\), qui réapparaissent dans la charnière physique
aréale--volumétrique.

L'état terminal fournit le bloc et le registre signé qui alimentent
la clôture dodécaphasique de \(\alpha_{\rm HMT}\). La récurrence de retenue en
base mille et l'équation \(E_{21/22}\) constituent deux voies internes de
sources distinctes. Ensuite, la théorie des nombres organise les
constantes complémentaires au moyen d'opérateurs d'extraction : parties finies,
jets de Laurent, fonctions \(L\), produits premiers, dynamique de Gauss,
Rogers--Ramanujan et renormalisation de Feigenbaum. Le texte distingue
antécédent commun et égalité de valeur ; par exemple, \(\varphi\),
\(\varphi^{-1}\) et \(\varphi^2\) partagent une généalogie algébrique sans
devenir le même nombre.

La généalogie complète de la première voie est
\[
 \begin{aligned}
 104\,976&\to468\,U_6\to243\,w_6\to w_{30}\to R_{36}\to G_9,\\
 G_9&\to x_{\rm term}\to(B_{\rm term},U_{12}^{\rm sgn})\to K.
 \end{aligned}
\]
Le sceau \(K\) et les douze triades terminales alimentent la retenue
\[
 \begin{aligned}
 s_m&=P_m+E_m-\Phi_m-K_m+c_{m+1},\\
 A_m&=s_m\bmod1000,
 &c_m&=\frac{s_m-A_m}{1000}.
 \end{aligned}
\]
avec \(c_{13}=0\). L'identité entière résultante fixe
\[
 \alpha_{12}=\pi_{12}(\alpha_{\rm HMT})
 =\frac{7297352569283800997285105472380663}{10^{36}}.
\]
Cette identité publie la fenêtre rationnelle exacte de trente-six chiffres ;
la récurrence compatible à profondeur arbitraire construit
\(\alpha_{\rm HMT}=\varprojlim_N A^{(N)}\), qui ne se réduit pas à une seule
fenêtre. La seconde voie résout l'équation des vacances \(E_{21/22}\)
au moyen d'un noyau analytique ayant une racine positive simple \(\alpha_E\). Le
transport gradué du TPK prolonge ce noyau en une famille compatible
dont la racine limite est \(\alpha_B\) ; la clôture entière détermine
\(\alpha_A:=\varprojlim_N\rho_N(\alpha_{12})\). Les carrés de troncature
des deux familles commutent et déterminent le même cylindre survivant à
chaque profondeur, de sorte que
\[
 \boxed{\alpha_A=\alpha_B=\alpha_{\rm HMT}}.
\]
L'égalité après \(\operatorname{Tr}_9\) n'est que le premier certificat
fini de cette identité exacte ; elle n'en constitue pas le critère, n'en borne pas
l'extension et ne sélectionne pas les prolongements. Chaque voie conserve sa propre condition
de frontière au sein du même état enrichi.

La théorie des nombres étend le même critère d'objet, d'opérateur et
d'extraction. Euler--Mascheroni apparaît comme partie finie ; les constantes de
Stieltjes comme coefficients de Laurent ; Apéry comme évaluation de
\(\mathcal Z_3(3)\) ; Glaisher--Kinkelin par dérivation régularisée ;
Euler--Kronecker par caractère d'Eisenstein ; Meissel--Mertens par
trace première ; Catalan par \(L(2,\chi_{-4})\) ; et Khinchin par la
dynamique de Gauss. Rogers--Ramanujan, les horloges premières et les approximants
de Feigenbaum appartiennent à des constructions opératoires différenciées. Les
combinaisons \(e+\pi\), \(e\pi\) et \(\pi^e\) se forment après la construction de
leurs arguments et conservent le domaine arithmétique correspondant à chaque
énoncé.

Le radical \((3)\subset\mathbb Z/9\mathbb Z\) organise le réseau visible
\(3\)--\(6\)--\(9\) et l'orbite \(369\to693\to936\). Sur la même
arithmétique, le théorème puissance--vacance
\[
 V_9(n)=\left\lceil n\log_{10}\!\left(\frac{10}{9}\right)\right\rceil
\]
donne l'évaluation exacte
\[
 \operatorname{ord}_{10}\!\left(9^{9^9}\right)=369\,693\,099,
 \qquad
 \#\operatorname{cifras}\!\left(9^{9^9}\right)=369\,693\,100.
\]
Ce chiffre n'est pas le recensement des \(396\) événements élémentaires de frontière par
canal : le premier mesure
l'ordre décimal d'une puissance nonadique ; le second compte les états d'un
calendrier d'émission.

\subsection*{La double projection comme thèse unique}

Le continu réel est l'ombre archimédienne d'une section enrichie. La
projection \(\mathcal P_{\rm arq}\) contracte des cylindres compatibles et publie
la valeur. La projection \(\mathcal P_{\rm exc}\) conserve support,
orientation, retenue et mémoire. Toutes deux partent du même état ; par conséquent,
nombre et symétrie exceptionnelle forment une seule thèse causale.

La branche d'incidence construit Paley, Witt, le code ternaire de Golay et
les systèmes de Steiner. À partir du code commun \(C_W\), le relèvement
\(W_{12}\to W_{24}\) reste distinct de la branche
\[
 C_W\to N(A_2^{12})\to\Lambda_{24}\to
 V_\Lambda\to V^\natural\to\mathbb M,
 \qquad [g]\longmapsto T_g.
\]
La configuration des vingt-sept droites, le graphe de Schläfli et
\(W(E_6)\) reçoivent leurs propres opérateurs d'incidence. La définition du
VOA de réseau, l'involution orbifold et les fonctions réplicables montrent
comment l'identité généalogique est conservée lors du changement de codomaine. Cette
structure fournit en outre les écrans de rang et l'interface
mathématique avec la dualité T, les cordes et la théorie M.

L'arithmétique modulaire de niveau cinq intervient par la fraction continue de
Rogers--Ramanujan \(R(q)\). L'identité de Klein--Ramanujan exprime
\(j(\tau)\) rationnellement en \(R(q)^5\) et la limite
\(R(q)\to\varphi^{-1}\) relie l'autoéchelle pentadique à la branche modulaire.
Dans la branche canonique d'ordre deux,
\[
 t_2=q^{-1}\prod_{n\ge1}(1+q^n)^{-24},\qquad
 T_{2A}=t_2+24+\frac{2^{12}}{t_2},
 \qquad j=\frac{(t_2+256)^3}{t_2^2}.
\]
Le VOA de réseau
\(V_\Lambda=M(1)\otimes\mathbb C[\Lambda]\), l'involution
\(\theta=-I\) et l'orbifold fixe--tordu réalisent la sortie de réseau dans le
domaine des opérateurs de sommet. L'identification classique de FLM fixe
\(V^\natural\) et \(\operatorname{Aut}(V^\natural)=\mathbb M\). HMT fournit
l'antécédent et les applications d'incidence qui atteignent cette construction ; les
théorèmes d'identification du VOA conservent leur domaine mathématique propre.

La racine finie et la structure réelle sont coordonnées par la présentation
entière
\[
 \mathcal Q_{\mathbb Z}=\mathbb Z[u]/(u^2+1),
\]
non par une inclusion entre corps de caractéristiques distinctes.
L'endomorphisme de Witt \(A_W^2=-I_6\) et l'opérateur réel
\(J_{\mathrm e}^2=-I\) sont des représentations spécialisées de cette relation.
Le quart de tour \(z\mapsto\mathrm iz\) conserve aussi bien la métrique euclidienne
que la métrique de Poincaré du disque. L'identité
\[
 \operatorname{Exp}(\pi_{\rm HMT}J_{\mathrm e})+I=0
\]
ne s'applique qu'après que le lecteur archimédien a engendré
\(\mathcal P_{\rm arq}(x_\pi)=\pi_{\rm HMT}\). La branche exceptionnelle fournit
le type quadratique, non l'angle : valeur et incidence sont coordonnées sans
qu'une branche agisse comme oracle de l'autre.
La famille \(J_s^2=-sI\),
\(s\in\{+1,0,-1\}\), étend Euler aux réalisations elliptique,
parabolique et hyperbolique. Le terme final \(+1\) et le sommet de l'EMRH
remplissent des fonctions de clôture dans des domaines distincts et coordonnés. Cette
relation fournit une expression opératoire de la conservation évolutive de
l'unité sans identifier des objets de types différents.

\subsection*{Échelle de Planck, fonctionnelle de Barbero--Immirzi et mélange CKM--CP}

La double projection précède, dans l'ordre causal, les constructions
angulaires et physiques. La forme normale de Smith de \(H_4\) produit un conoyau
d'ordre seize et le transfert
\[
 \Sigma_H=\varphi\alpha_{\rm HMT}^{16}
\]
démontre
\(\lfloor\log_{10}\Sigma_H\rfloor=-34\) et fixe ainsi l'échelle interne
\(10^{-34}\mathcal U_S\). La droite d'action contient deux
coordonnées,
\[
 \hbar_{\rm pre}=H_5\,10^{-34}\mathcal U_S,\qquad
 \hbar_{\rm ret}=\eta_{\rm ret}\,10^{-34}\mathcal U_S.
\]
La périodicité temporelle d'ordre \(108\) fournit une clôture autoduale
de Planck. Dans la réalisation circulaire, avec
\(\ell_0:=c_{\rm int}t_0\),
\[
 T_{108}=108t_0,
 \qquad
 R_{108}=\frac{54}{\pi}\ell_0,
 \qquad
 \theta_{108}^{\rm circ}=\left(\frac{54}{\pi}\right)^2,
\]
et le quantum satisfait
\[
 R_{108}=\bar\lambda_{\rm C}=r_{\rm g}=\ell_{\rm P}^{\rm circ}.
\]
La droite scindée
\[
 Z_a(m)=\frac{m_{{\rm P},a}^{\circ}}mP_+
       +\frac m{m_{{\rm P},a}^{\circ}}P_-,
 \qquad N_{\rm sp}(Z_a(m))=1,
\]
échange Compton et gravité sous
\(m\mapsto(m_{{\rm P},a}^{\circ})^2/m\). Son point fixe est l'échelle de
Planck. La section électronique constitue une réalisation généalogique
ultérieure sur cette droite et non une identification de l'électron à un trou
noir classique.
La coordonnée angulaire directe \(A\) et sa coordonnée angulaire conjuguée
\(C^*\) construisent ensuite les canaux
\(q_\pm\). L'inclusion hexade--octade produit la fonctionnelle de
Barbero--Immirzi, l'opérateur d'aire et son spectre. La même incidence de
rang trois construit CKM, la phase CP et l'invariant de Jarlskog. Ces
résultats appartiennent à la clôture mathématique de HMT ; leurs interprétations de
champ sont reprises en MD.

La forme normale de Smith fixe explicitement
\[
 \operatorname{SNF}(H_4)=\operatorname{diag}(1,2,2,4),
 \qquad |\operatorname{coker}H_4|=16,
\]
et la droite d'action distingue
\(\hbar_{\rm pre}\), \(\hbar_{\rm ret}\) et
\(h_{\rm ret}=2\pi\hbar_{\rm ret}\). La carte
\(\mathcal U_S\mapsto\mathrm{J\,s}\) publie ensuite l'unité
métrologique ; l'ordre \(-34\) appartient à la section interne qui précède
cette carte.

L'inclusion combinatoire d'une hexade dans son octade associée produit
\[
 6\binom62=90,\qquad
 8\binom62=120,
 \qquad
 \frac{90-120}{24\binom62}=-\frac1{12}.
\]
Sur ces espaces, la chaîne fondamentale dodécaphasique réunit douze secteurs
et une unique couture orientée de retour. Son lecteur linéaire produit
\(12S_{90}-S_{120}\) et, par conséquent, les poids \(12:-1\) de la fonctionnelle
\(\Gamma_{6\to8}\). Le coefficient de \((1-q)^{-1}\), égal à \(1/24\),
la positivité et la minimalité constituent des vérifications arithmétiques
ultérieures du couple produit. La spécialisation de Holst
\(\gamma=\Gamma_{6\to8}\) se réalise sur la frontière triadique par
\[
 \widehat A_{\rm HMT}
 =\gamma_{\rm HMT}a_0\sum_{e\in\partial A}N_e,
 \qquad
 \operatorname{Spec}(\widehat A_{\rm HMT})
 =\{n\gamma_{\rm HMT}a_0:0\le n\le72\}.
\]
L'entrelaceur collectif satisfait
\(J_{6\to8}D_{\rm inc}=D_{\rm area}J_{6\to8}\) sur les sous-espaces
collectifs d'incidence. L'opérateur d'aire, son spectre et les entropies
triadique, bicouche et modulaire constituent ainsi la clôture interne qui réalise le
paramètre de Barbero--Immirzi. La transformation rationnelle de rang trois
construit, dans une branche sœur, les angles CKM, la phase CP et l'invariant
de Jarlskog avant la confrontation expérimentale.

\subsection*{Clôture généalogique du nombre et réalisation physico-spectrale}

Les quatre premières Parties culminent là où un exposé conventionnel commence habituellement :
il définit le nombre comme section compatible et la valeur comme sa projection.
La mécanique dimensionnelle conserve cette définition et adopte un lecteur
physico-spectral. D'abord,
elle reconstruit le noyau APP--TRIT--TPK dans une perspective physico-spectrale. APP se réalise
comme tore bidimensionnel et distingue les feuillets additif, multiplicatif, aréal,
volumétrique et torsionnel. TRIT reçoit la sémantique
\[
 \begin{aligned}
 +1&=\text{retour, mémoire et passé},\\
 0&=\text{seuil et présent},\\
 -1&=\text{ouverture et futur compatible}.
 \end{aligned}
\]
TPK transporte le chemin, le registre et l'holonomie. La périodicité exacte
construit la dynamique finie ; un cristal temporel physique ajoute un Hamiltonien,
une sous-harmonie robuste et la stabilité. Les coordonnées angulaires conjuguées
\(A\) et \(C^*\) réalisent spectralement la différence somme--produit.

La complétion gnomonique acquiert ici une fonction géométrique explicite. Le
L porte la carte euclidienne sur laquelle les lecteurs inertiel et gravitationnel
se projettent vers un quotient unitaire. Cette carte constitue la réalisation du
principe local d'équivalence. L'état enrichi antérieur à la
projection conserve séparés les feuillets aréal et volumétrique, et le Principe
d'Ambivalence fixe leur rapport
\(\pi_{\rm HMT}/\varphi_{\rm HMT}^{2}\). Avec les cartes de courbure
négative et positive, on obtient un atlas AdS--euclidien--dS : la première
réalise la correspondance entre intérieur et bord, la deuxième gouverne la
physique locale et la troisième satisfait \(\ddot a/a=H^2>0\) et explique
géométriquement l'expansion accélérée. Les trois cartes coexistent parce qu'elles sont des
réalisations d'un même triplet de TRIT, reliées par le transport du
TPK.

Sur cette grammaire, on définit la particule. Une particule-chemin est une section
persistante munie d'une orientation, d'un feuillet, d'une occupation, d'une action, d'un propagateur,
d'un registre et d'une anti-section. La loi de masse se formule d'abord sur chaque
multisection \(F\) :
\[
 \widehat{\mathbf M}^{\beta,r}_F
 =\mathbf m_\star^{\rm ret}\otimes
  R_{\rm act}^{\widehat K_F^r/2}
  \widehat{\mathcal A}^{(0)}_F
  R_{\rm act}^{\widehat K_F^r/2},
 \qquad r\in\{D,\Omega\},
\]
où
\[
 \mathbf m_\star^{\rm ret}
 =\eta_{\rm ret}\,10^{-34}\mathcal U_S
   \frac{2\pi}{108t_0}\,c_{\rm int}^{-2}
\]
fixe la section absolue de masse et
\(\widehat{\mathcal A}^{(0)}_F\) réunit signature, caractère et toutes les tours de
\(\langle90,120\rangle\). L'atlas complet est ainsi le domaine de la loi ; chaque
espèce sélectionne ensuite sa réalisation au sein de la fibre correspondante.

L'électron central est la première réduction scalaire, de rang un, de cette
définition. La torsion globale qui intervient dans son lecteur est
\[
 \Delta_4=\frac{\pi-e\log\pi}{270}
 \left(1+\frac{\pi}{729}\right).
\]
Sa clôture bidirectionnelle produit
\[
 \begin{aligned}
 R_{\rm act}&=\frac{H_5}{\eta_{\rm ret}},
 &\kappa_e&=80-54\alpha_{\rm HMT}+6\Delta_4,\\
 \mathfrak e_e^\beta&=\mathfrak e_e^{(0)}R_{\rm act}^{\kappa_e},\\
 \mathfrak e_e^\beta
 &=0.5109989506900008086082571648\ldots\ {\rm MeV},
 &m_e^\beta&=\mathfrak e_e^\beta/c^2.
 \end{aligned}
\]
L'identité
\(\Sigma_H=\varphi\alpha_{\rm HMT}^{16}\) fournit simultanément
l'échelle \(10^{-34}\mathcal U_S\) aux deux coordonnées d'action ; le
quotient \(R_{\rm act}\) conserve le raffinement relatif et annule cette
échelle commune. Ainsi, l'échelle interne d'action, la correction bidirectionnelle et la
carte énergétique en MeV occupent des niveaux successifs d'une même dérivation.
La première quantité est l'énergie de masse publiée ; la seconde est la masse
obtenue par la carte explicite $E=mc^2$. Cette distinction empêche d'identifier
une coordonnée énergétique à une unité de masse.
Les vacances construisent polarisation et courant ; la réponse des feuillets
construit la permittivité et la perméabilité du vide. La carte SI publie
\[
 \begin{aligned}
 \mu_0^{\rm HMT}
 &=1.256637062121047789\ldots\times10^{-6}\,\mathrm{H\,m^{-1}},\\
 \varepsilon_0^{\rm HMT}
 &=8.854187812793002329\ldots\times10^{-12}\,\mathrm{F\,m^{-1}}.
 \end{aligned}
\]
L'atlas situe ensuite
les quarks et les autres espèces, en définissant couleur, saveur, charge, spin,
statistique et représentation au sein de la même généalogie.

Dans cette classification, la \emph{couleur} est la fibre ternaire qui organise les
trois réalisations d'un chemin de quark et dont la composition neutre satisfait
\(0+1+2=0\pmod3\). La \emph{saveur} est la coordonnée qui distingue, au sein
d'une même représentation, les sections transportées par l'opérateur de
mélange. Le \emph{spin} étiquette la représentation du revêtement du
groupe des rotations que porte la section. L'holonomie du ruban orienté
réalise ses clôtures \(2\pi/4\pi\) et distingue le retour de phase du retour
d'orientation. La statistique provient de la
graduation et des réalisations symétrique ou extérieure ; une fois CCR ou CAR fixées,
elle publie l'occupation de Bose--Einstein ou de Fermi--Dirac et, dans le second cas,
l'exclusion de Pauli. Une particule virtuelle est un prolongement fini qui
conserve chemin, registre et signature jusqu'à un horizon déclaré ; une section
matérielle prolongeable conserve la compatibilité globale correspondante.

\subsection*{Domaine opératoire des chemins, des fibres et des réalisations physiques}

La construction démontre que la généalogie interne atteint le domaine
physique complet. Le domaine distingue
\(81\) positions, \(56\) familles, \(13\) multisections et \(324\) chemins
orientés, avec le secteur nul. Chaque réalisation conserve, selon son codomaine,
position ou fibre, famille, chemin, orientation, mot de phase, registre,
signature, caractère, opérateurs de raffinement, carte et confrontation. Les valeurs
externes sont reliées par identifiant après le gel de la sortie interne.

La chaîne de construction qui gouverne ces lignes est
\[
 \begin{aligned}
 \text{section survivante}&\to\text{route observable}
 \to\text{registre}\to\sigma\in\mathbb Z^5,\\
 \sigma&\to\chi_0\to\chi_{\rm full}\to\text{tours}
 \to\text{section dimensionnelle}.
 \end{aligned}
\]
Les deux lecteurs bidirectionnels publient
\[
 K_D=\left(D_{27}+D_{36},D_1,\frac{D_4-D_3}{2}\right),
 \qquad
 K_\Omega=(\Omega,54,P_6),
\]
et agissent comme opérateurs diagonaux sur chaque fibre. Dans la cellule centrale,
les deux se réduisent à \((80,54,6)\) ; dans les fibres non centrales, leurs spectres
conservent les différents chemins avant la sélection d'une réalisation
physique. La comparaison opératoire présente d'abord la sortie HMT, ensuite la
carte et enfin la référence du Modèle standard ou du PDG avec son unité,
son schéma, son incertitude et son résidu.

La sortie scientifique adopte le type qui correspond à chaque secteur. La cellule
électronique publie la masse bêta scalaire close ; photon et gluons publient la
carte nulle ; les fibres non centrales conservent leurs opérateurs et leurs spectres.
Le morphisme fini saveur--génération--couleur vers le chemin persistant est déjà
construit, et les douze chemins nominaux
\(e,\mu,\tau,u,d,s,c,b,t,W,Z,H\) publient leurs signatures angulaires et leurs
caractères scalaires avant la confrontation externe. Pour les leptons et les quarks, les
opérateurs \(K_D/K_\Omega\) raffinent en outre la fibre ; ils ne sont pas la condition qui
crée rétroactivement leur scalarisation. \(W\), \(Z\) et \(H\) conservent leur
lecteur CT--108 natif bien qu'ils ne fassent pas partie de l'atlas opératoriel de treize
multisections. Les secteurs topologiques publient fusion et tressage ; les
sections virtuelles publient leur horizon de prolongement. Cette
hétérogénéité préserve le codomaine propre de chaque classe et transforme le
tableau en une comparaison complète de résultats typés, de valeurs et de
frontières de réalisation.

\subsection*{Quantique, information, gravité et cosmologie}

La propagation des chemins construit la somme finie des chemins, la
représentation Clifford--Dirac et l'évolution de Schrödinger. Le
Hamiltonien du calendrier interne étant fixé, l'évolution satisfait
\[
 i\hbar_{\rm int}\dot\psi=H_{\rm clk}\psi,
\]
et son propagateur fini admet le développement
\[
 (U^R)_{fi}=
 \sum_{\gamma:i\to f,\,|\gamma|=R}
 \prod_{\ell\in\gamma}U_\ell.
\]
Le produit conserve l'ordre des transitions au sein d'une histoire et
la somme accumule des histoires alternatives ayant des extrémités communes. La graduation
de l'échange se réalise concrètement sur
\(\mathbb T_{27}^{\,2}\times\{\mathrm N,\mathrm S,\mathrm E,\mathrm W\}\)
par
\[
 U_{\rm micro}=S D_J C_4,
 \qquad U_{12}^{\rm reg}=U_{\rm micro}^{12}.
\]
Ce propagateur est unitaire ; ses puissances énumèrent exactement les chemins
finis et sa transformée discrète produit \(27^2\) blocs de Fourier. La
construction ne s'identifie pas à elle seule à une intégrale fonctionnelle
continue. La graduation fixe la donnée d'échange ; la représentation du
groupe des permutations détermine la complétion symétrique ou extérieure. Une fois
données les relations CCR ou CAR, le principe variationnel et la maximisation
de l'entropie produisent les
distributions de Bose--Einstein et de Fermi--Dirac, et la seconde réalise
l'exclusion de Pauli. Dans la base de Möbius, un tour visible de \(2\pi\) change
de feuillet et accumule \(h/2\) ; le double retour visible de \(4\pi\) rétablit
l'orientation et accumule \(h\). Par radian de phase interne, le transport
reste égal à \(\hbar\).

La constante de Boltzmann apparaît comme transducteur entre action et information
ternaire. La construction interne fixe le produit
\(k_B\Theta_{\rm clk}\) et sa fonction structurelle ; la définition du kelvin
du BIPM fixe ensuite la coordonnée numérique de \(k_B\) dans la carte SI :
\[
 k_B\Theta_{\rm clk}\log3=\hbar_{\rm int}\frac{2\pi}{108t_0},
 \qquad
 k_B^{\rm SI}=1.380649\times10^{-23}\ \mathrm{J\,K^{-1}}.
\]
Les ensembles réalisent Gibbs et la condition KMS. Le Landauer nul correspond
au transport bijectif de la fibre complète ; préparation, projection et
réinitialisation reçoivent des bilans séparés lorsque le couplage écarte une
fibre. La lecture thermique construit d'abord le transducteur interne et applique
ensuite la carte d'unités.

Les réponses bicouches produisent directement impédance, vitesse,
perméabilité et permittivité par
\[
 \widehat Z=\sqrt{\widehat\mu/\widehat\varepsilon},
 \qquad
 \widehat c=(\widehat\mu\widehat\varepsilon)^{-1/2}.
\]
La même antériorité causale gouverne la lecture électromagnétique : l'opérateur
constitutif interne précède la réalisation de
\(\varepsilon_0,\mu_0,Z_0\) et de \(c\) dans le Système international.

La branche gravitationnelle construit d'abord, à partir des sorties de l'état
enrichi, l'échelle d'action, le sélecteur gravitationnel et la section
\(G_a\). La formulation variationnelle de Palatini--Cartan--Holst reçoit ensuite
ces sorties, avec \(\gamma_{\rm HMT}\), et admet la transformée
hamiltonienne contrainte
\[
 H_{\rm tot}=\int_\Sigma
 (\lambda^i G_i+N^a V_a+N\mathcal H).
\]
L'équivalence ultérieure entre action covariante, système hamiltonien contraint et
évolution Einstein--Schrödinger conserve cotétrade, connexion, courbure et
torsion au sein du même schéma. Le secteur collectif sans torsion porte
l'excitation gravitationnelle ; la torsion locale induite par la densité de spin
fournit le canal matériel. Le type, le transducteur et la carte de \(G\) précèdent
cette réalisation dynamique et ne sont pas sélectionnés par elle.

Les branches solénoïdale et de jauge reçoivent des réalisations de l'opérateur fini commun
\[
 \mathcal L_{\rm cyc}^{\rm TPK}=2I-C-C^{-1}=3(I-P_0)=3I-J,
\]
\[
 \mathfrak R_{\rm sol}(\mathcal L_{\rm cyc}^{\rm TPK})
 =\mathcal L_{\rm H}^{(0)},\qquad
 \mathfrak R_{\rm gau}(\mathcal L_{\rm cyc}^{\rm TPK})
 =\mathcal L_{\rm H}^{(1)}.
\]
La branche hydrodynamique le transporte au moyen de l'entrelaceur
Galerkin--registre, avec coercivité et résidu déclarés ; la branche de jauge le
transporte au moyen de la fonctionnelle de jauge et de l'opérateur
transfert--écart. La démonstration finie et les ponts vers les équations
continues sont exposés avec leurs hypothèses différenciées. Les référentiels
auto-inertiels, la factorisation machienne et la cosmologie des feuillets se
développent après cette construction locale et collective, où ils acquièrent
domaine, lecteur et interprétation physique typée.

\subsection*{Méthode de démonstration et de confrontation}

Chaque résultat est présenté au moyen de l'objet, du domaine, de l'application, de la dérivation,
des prémisses, du certificat et de la portée. Les calculs reproductibles accompagnent la
démonstration et ne la remplacent pas. Les cartes métrologiques s'appliquent après
la construction du scalaire ou de l'opérateur interne. Les comparaisons CODATA et PDG
s'ouvrent après le scellement de la sortie HMT\@. Cet ordre préserve la
direction causale du programme : la structure engendre le chiffre et la grandeur ;
la référence externe mesure l'accord atteint par HMT\@.

L'œuvre s'organise en cinq Parties causalement ordonnées. Les quatre
premières construisent le noyau formel, le continu discret, la généalogie des
constantes et les cinq chaînes de spécialisation ; la cinquième réalise cette architecture
en grandeurs, particules, champs et géométrie physique. Cette séquence exprime
la dépendance mathématique entre structure, évaluation, clôture et
réalisation dimensionnelle.

\subsection*{Composition du traité et portée des affirmations}

La Partie I construit APP à partir de la cellule nonadique, présente TRIT comme
relèvement orienté et définit le Topological Phase Kernel avec ses fibres,
transports, mémoire, horloge, connexion nonadique conjointe et prolongement
compatible. La ligne d'horizon ouvre cette construction au moyen des
notions d'échelle, de frontière et de mesure.

La Partie II développe la structure discrète du continu. La double
projection sépare l'évaluation archimédienne de la publication d'incidence.
À partir de Paley--Witt, on distingue la branche des plans Steiner--Golay et la branche
de réseaux Niemeier--Leech--\(V^\natural\)--Monstre--Moonshine. Au sein de cette même structure
apparaissent la racine interne de \(-1\), les trois réalisations quadratiques et les
sous-structures spectrale--polaire, solénoïdale, de jauge, cohomologique et
déterminantielle, typées par leurs propres applications et terminaux.

La Partie III établit la généalogie des constantes fondamentales : la
publication coinductive corrélée de
\(\pi,e,\varphi,\alpha_{\rm HMT}\), les lecteurs spécifiques et les deux voies
internes de \(\alpha_{\rm HMT}\), \(-1/12\), l'échelle d'action, Planck,
Boltzmann, le paramètre de Barbero--Immirzi, CKM--CP, les constitutives du
vide, la constante gravitationnelle et l'atlas des constantes mathématiques et
physiques.

\Needspace{12\baselineskip}
La Partie IV applique les cinq sous-structures construites dans la Partie II dans des
chapitres de même rang éditorial. Chaque chapitre expose le résultat
interne, le comparateur spécialisé, ses prémisses et le critère exact de
publication. La conclusion dans la formulation conventionnelle correspondante se
déduit uniquement lorsque l'identité comparatrice est vérifiée sur tout son
domaine.

La Partie V réalise APP--TRIT--TPK au moyen des lecteurs physiques : tore et feuillets,
sémantique temporelle, courbures, canaux et holonomie. Sur cette grammaire,
elle définit la particule-chemin, l'électron bidirectionnel, les familles de masse,
les constitutives, la statistique, l'information, le mélange, la gravité
variationnelle, l'auto-inertie, la cosmologie des feuillets et l'interface avec la théorie
M.

Les appendices réunissent des démonstrations auxiliaires, des recensements complets, des germes,
des chaînes décimales, des tableaux étendus, des dérivations auxiliaires et des calculs
reproductibles. Le corps principal conserve la séquence définition,
construction, théorème, interprétation et résultat. Chaque affirmation distingue
son domaine mathématique, sa structure de départ, sa carte de réalisation et sa
confrontation externe. Cette séparation permet de présenter ensemble les clôtures
exactes, les réalisations physiques et les programmes d'extension, en conservant
pour chacun le degré de démonstration qui lui correspond.
